% ============================================================================
% Stitched Supplementary Information
% Generated by merge_si.sh on 2026-06-12T18:10:37Z
% Section order: generative → PES → kinetics → shared → seeds → parameters
% ============================================================================

% --- begin si2_generative_loop.tex ---
% ============================================================================
% Supplementary Information SI-2: Generative-Loop Reaction Discovery
%
% Drop-in subsection fragment. Parent SI must provide: amsmath/amssymb,
% booktabs, siunitx (with \angstrom). Bibliography handled at the parent
% level. (Parameter reference lives in si_parameters.tex.)
% ============================================================================

\subsection{Generative-Loop Reaction Discovery}
\label{si:generative}

Each iteration processes a batch of $32$ \textbf{exploration contexts} (a reactant geometry --- single molecule or merged pair --- plus the artefacts the pipeline attaches downstream: TS conformer, product conformer, IRC trajectories, fragment list, barriers). The pipeline has eight stages: (1)~sampling, (2)~merge enhancement, (3)~MoreRed diffusion, (4)~forward barrier gate, (5)~Hessian + IRC validation, (6)~endpoint identity, (7)~fragmentation detection, (8)~backward barrier gate. Critic stages (4--8) reuse the shared machinery in SI-\ref{si:shared} (ML force field, geometry primitives, optimisers); this section documents what is specific to the generative path.


%-----------------------------------------------------------------------------
\subsubsection{Stage 1 --- Concentration-weighted context sampling}
\label{sec:sampling}
%-----------------------------------------------------------------------------

Half of each batch is single-molecule, half is bimolecular merge pairs. The eligible compound pool is scoped to the current experiment and excludes out-of-scope frontier compounds.

\paragraph{Single-molecule sampling.}
A species is drawn from the kinetic snapshot's steady-state distribution (SI-\ref{si:kinetics}); compounds with $n_\text{atoms} < 3$ are excluded. Within the chosen compound, the conformer is drawn from the room-temperature Boltzmann distribution,
\begin{equation}
    P(m) = \frac{\exp(-(E_m - E_\text{min}) / k_B T)}{\sum_{m'} \exp(-(E_{m'} - E_\text{min}) / k_B T)}, \qquad k_B T = \SI{0.0257}{eV}.
\end{equation}

\paragraph{Bimolecular pair sampling.}
$50 \times n_\text{needed}$ candidate pairs are drawn from the steady-state distribution, then filtered against (i)~a composition cap of $\{C{:}5,\ H{:}10,\ O{:}5\}$ (tetrose-sized merged complexes) and (ii)~a per-pair coverage cap of $4$ non-manual reactions in the current experiment. For each surviving pair, one minimum is drawn from each compound at random. Duplicate pairs are kept --- each becomes an independent rotation/noise seed downstream.

\paragraph{Initial geometric merge.}
Each component receives an independent random $SO(3)$ rotation; a random unit vector defines the separation direction; the smallest displacement that yields a non-negative minimum van der Waals gap $\min_{ij}\bigl(\lVert r_i - r_j \rVert - r^{\text{vdW}}_i - r^{\text{vdW}}_j\bigr)$ is found by binary search, then $\SI{0.2}{\angstrom}$ of repulsive buffer is added. Atoms are interleaved by stable-sorting on atomic number; the merged structure is COM-centred. Atom ordering is fully deterministic so that downstream IDs match across workers.

\noindent\begin{tabular}{lr}
\toprule
Parameter & Value \\
\midrule
Batch size & 32 contexts \\
Single/merge split & 50/50 \\
Oversample factor (merges) & 50 \\
Min compound size & 3 atoms \\
Composition cap (merge) & $\{C{:}5,H{:}10,O{:}5\}$ \\
Per-pair reaction cap & 4 \\
$k_B T$ (conformer pick) & \SI{0.0257}{eV} \\
vdW separation buffer & \SI{0.2}{\angstrom} \\
\bottomrule
\end{tabular}

The 50/50 split is hardcoded because a uniform pair draw with $n$ compounds spends $\sim 1/n$ on each cross-pair, making bimolecular discovery rare without a fixed budget. The tetrose composition cap is the dominant top-down constraint on system reach; relaxing it is the principal scale-up knob for chemistry beyond formose-style sugars.


%-----------------------------------------------------------------------------
\subsubsection{Stage 2 --- Merge encounter enhancement}
\label{sec:merge-prep}
%-----------------------------------------------------------------------------

The random encounter geometry from Stage 1 is geometrically valid but typically far from any saddle. Pre-conditioning reduces the burden on the diffusion stage. For each merge context: (i)~generate $5$ additional random merge geometries and pick the lowest-ML-energy candidate; (ii)~run $5$ FIRE micro-relaxation steps with $f_\text{max} = \SI{0.05}{eV/\angstrom}$ to settle the fragments into a reasonable encounter geometry without converging to the encounter-complex minimum; (iii)~reject the context if the binding energy $E_\text{bind} = E(A) + E(B) - E(\text{complex}) > \SI{0.3}{eV}$ (deep encounter complexes are kinetic traps from which both IRC directions relax back into the complex); (iv)~apply per-atom Gaussian noise of $\sigma = \SI{0.15}{\angstrom}$ to break symmetry, since MoreRed is equivariant and produces no movement on a symmetric stationary input; (v)~re-centre for the translation-invariant diffusion.

\noindent\begin{tabular}{lr}
\toprule
Parameter & Value \\
\midrule
Extra orientations ($K$) & 5 \\
FIRE micro-relax steps & 5 \\
FIRE $f_\text{max}$ & \SI{0.05}{eV/\angstrom} \\
Symmetry-break $\sigma$ & \SI{0.15}{\angstrom} \\
Encounter-complex cap & \SI{0.3}{eV} \\
Min pairwise distance (overlap guard) & \SI{0.8}{\angstrom} \\
\bottomrule
\end{tabular}

The \SI{0.3}{eV} binding cap is a chemistry choice: physisorbed and H-bonded complexes are kept, while covalently bound encounter complexes (themselves minima) are rejected --- their reactions out are PES-loop business (SI-\ref{si:pes}). The \SI{0.15}{\angstrom} displacement is calibrated against the diffusion-model noise scale: at $t \sim 200$ on the polynomial schedule the cumulative variance is comparable, so the denoiser sees a perturbation it can act on without losing the input identity.


%-----------------------------------------------------------------------------
\subsubsection{Stage 3 --- Diffusion-based TS proposer}
\label{sec:diffusion}
%-----------------------------------------------------------------------------

\paragraph{Model.}
A SchNetPack equivariant graph network (Sch{\"u}tt et al.\ 2023), trained as a MoreRed-JT joint-time denoiser (Kahouli et al.\ 2024) on relaxed transition states. Per atom it predicts a noise vector $\boldsymbol{\epsilon}_\theta(\mathbf{x}_t, t) \in \mathbb{R}^3$ and a time-step $\hat{t}_\theta(\mathbf{x}_t)$ used for adaptive convergence. The forward process is a variance-preserving Gaussian DDPM (Ho et al.\ 2020) of length $T = 1000$ with a polynomial noise schedule (Hoogeboom et al.\ 2022),
\begin{equation}
    \bar{\alpha}_t = (1 - s)\bigl(1 - (t/T)^p\bigr)^2 + s, \qquad p = 2,\ s = 10^{-5},
\end{equation}
clipped so $\alpha_t / \alpha_{t-1} \geq 10^{-3}$, with forward conditional $q(\mathbf{x}_t \mid \mathbf{x}_0) = \mathcal{N}\bigl(\sqrt{\bar{\alpha}_t}\,\mathbf{x}_0,\,(1 - \bar{\alpha}_t)\,\mathbf{I}\bigr)$ and the standard DDPM lower-bound posterior variance for the reverse step. Translation invariance is enforced by re-centring at every step.

\paragraph{Partial-noise inference.}
Unlike text-to-image diffusion, the TS proposer does not sample from pure Gaussian noise: it partially noises a chemically meaningful input and then denoises. Each batch samples a single noise level $t \sim \mathcal{U}[100, 650]$, forward-diffuses each starting geometry as $\mathbf{x}_t = \sqrt{\bar{\alpha}_t}\,\mathbf{x}_0 + \sqrt{1 - \bar{\alpha}_t}\,\boldsymbol{\epsilon}$ with $\boldsymbol{\epsilon} \sim \mathcal{N}(0, I)$, and reverse-diffuses with up to $1000$ steps:
\begin{equation}
    \mathbf{x}_{t-1} = \tfrac{1}{\sqrt{\alpha_t}}\bigl(\mathbf{x}_t - \tfrac{1-\alpha_t}{\sqrt{1 - \bar{\alpha}_t}}\,\boldsymbol{\epsilon}_\theta\bigr) + \sigma_t\,\mathbf{z}, \quad \mathbf{z} \sim \mathcal{N}(0,I).
\end{equation}
Because the model predicts a per-molecule time-step, the reverse loop is adaptive: molecules whose internal clock decays faster converge in fewer steps, while harder cases use more. Once a molecule's predicted time reaches the convergence threshold ($0$) it is held fixed while the rest continue.

\paragraph{Outputs.}
Each context receives a denoised \texttt{ts\_conformer} (COM-centred), an ML \texttt{ts\_energy}, a geometry hash for cross-worker deduplication, and discovery metadata (\texttt{discovery\_method = "generative"}, the injected noise level $t$, timestamp).

\noindent\begin{tabular}{lr}
\toprule
Parameter & Value \\
\midrule
$T$ (schedule length) & 1000 \\
Schedule shape & polynomial, $p=2$, $s = 10^{-5}$ \\
Clip floor & $10^{-3}$ \\
Variance type & lower-bound \\
Partial-noise range & $[100, 650]$ \\
Max reverse steps & 1000 \\
Convergence step & 0 \\
Cutoff (neighbour list) & \SI{5.0}{\angstrom} \\
\bottomrule
\end{tabular}

The $[100, 650]$ noise range is the central hyperparameter: too low and the denoiser barely moves the input (most candidates are reactant copies); too high and the input is washed out (chemistry is lost). The chosen range keeps $\bar{\alpha}_t \in [0.86, 0.005]$, the regime where the noised input still encodes meaningful structure but the denoiser has freedom to move atoms across bonds. The injected $t$ and the model-predicted $\hat{t}_\theta$ are not constrained to agree --- the reverse loop adapts to the noise level it actually perceives, which gives robustness to off-distribution inputs (e.g., symmetry-broken encounter complexes from Stage 2).


%-----------------------------------------------------------------------------
\subsubsection{Stage 4 --- Forward barrier gate}
\label{sec:fwd-barrier}
%-----------------------------------------------------------------------------

A cheap ML-energy filter. The forward barrier is
\begin{equation}
    \Delta E_\text{fwd} =
    \begin{cases}
        E_\text{TS} - (E_A + E_B) & \text{merged, components } A,B,\\
        E_\text{TS} - E_\text{start}  & \text{single-molecule,}
    \end{cases}
\end{equation}
accepted iff $\Delta E_\text{fwd} \in [\SI{0}{eV}, \SI{2.7211}{eV}]$ (the upper bound is $0.1$~Hartree).

The \SI{2.7211}{eV} ceiling is deliberately loose. At \SI{300}{K} the Eyring prefactor $k_B T/h \approx 6 \times 10^{12}\ \text{s}^{-1}$, so even \SI{2.7}{eV} corresponds to a unimolecular rate $\sim 10^{-33}\ \text{s}^{-1}$ --- kinetic relevance is enforced downstream by the ODE solver, not here. Tightening this cap would risk discarding reactions whose ML barrier is overestimated relative to DFT.


%-----------------------------------------------------------------------------
\subsubsection{Stage 5 --- Hessian and IRC validation}
\label{sec:hessian-irc}
%-----------------------------------------------------------------------------

The bidirectional IRC algorithm (Hessian projection, adaptive displacement, bidirectional Newton relax, energy-based direction assignment) is shared with the PES loop and described in \S\ref{sec:shared-irc}. Generative-specific choices:

\paragraph{Hessian acceptance.}
The candidate is accepted as a TS iff \emph{at least one} significant eigenvalue ($|\lambda_i| > 10^{-4}\,\text{eV/\AA}^2$; \S\ref{sec:shared-hessian}) satisfies $\lambda_i < \SI{-0.01}{eV/\angstrom^2}$. The eigenvector of the \emph{most} negative eigenvalue is taken as the imaginary mode, and this single mode is followed by the bidirectional IRC. We deliberately do \emph{not} require the strict ``exactly one negative eigenvalue'' criterion that the PES loop enforces (\S\ref{sec:pes-prfo}): higher-order signatures are common at diffusion-proposed geometries and forcing strict first-order acceptance would tank recall.

\emph{Consequences of admitting higher-order saddles.} Once the candidate passes this gate, no downstream stage re-counts the negative eigenvalues. A geometry with $\geq 2$ negative modes therefore enters the pipeline if the dominant mode is selected; the subsequent IRC, endpoint, and barrier gates filter such candidates only \emph{implicitly}: a true higher-order saddle is statistically unlikely to (i) relax cleanly to two distinct minima within $350$ Newton steps on each side, (ii) have its dominant mode connect two endpoints that satisfy the endpoint-identity test (Stage 6), and (iii) yield both a forward and a backward barrier inside $(0, \SI{2.7211}{eV}]$ (Stages 4, 8). Pathological higher-order candidates typically fail one of these implicit filters --- most often the bidirectional-barrier check, where the secondary negative mode causes one endpoint to land at an energy above the TS. But candidates whose dominant mode is a sensible reaction coordinate \emph{can} survive even when secondary negative modes are present; we accept that this means the generative loop occasionally admits geometries that are not strictly first-order saddles. Transition-state theory assumes a first-order saddle for the Eyring prefactor, so these admissions carry a small kinetic-rate bias, mitigated downstream by the PBE0/def2-TZVPP barrier refinement (\S\ref{sec:kinetics-dft}) which is computed at the same in-box geometry and re-anchors the rate constants on DFT energetics.

\paragraph{Adaptive displacement constants.}
The clip range is wider than the PES loop's because diffusion-proposed saddles span a broader curvature distribution:
\begin{equation}
    \Delta x = \mathrm{clip}\bigl(\tfrac{0.2}{\sqrt{|\lambda_\text{TS}|}},\ \SI{0.05}{\angstrom},\ \SI{0.5}{\angstrom}\bigr).
\end{equation}

\paragraph{Endpoint relaxation budget.}
Trust-region Newton (\S\ref{sec:shared-newton}) with max $350$ steps (longer than the PES loop's $120$, again reflecting the broader candidate distribution), $f_\text{max} = \SI{0.005}{eV/\angstrom}$, Hessian re-traced every $5$ steps. Acceptance: convergence or bail-late at $f_\text{max} < 10\times$ tolerance.

\paragraph{Context rebuild after direction flip.}
When the IRC reactant (energy-determined; \S\ref{sec:shared-irc}) differs from the sampler's source compound, the context is rebuilt against the actual reactant: re-fragment, re-derive SMILES, re-relax each fragment on its own PES, re-derive merge/single status. This is what lets the generative loop discover reactions whose reactant the sampler did not anticipate.

\noindent\begin{tabular}{lr}
\toprule
Parameter & Value \\
\midrule
Eigenvalue threshold (TS) & \SI{-0.01}{eV/\angstrom^2} \\
Significance mask & $10^{-4}\,\text{eV/\AA}^2$ \\
Adaptive step formula & $0.2 \cdot |\lambda|^{-1/2}$~\AA, clamped $[0.05, 0.5]$ \\
IRC max steps / direction & 350 \\
IRC $f_\text{max}$ & \SI{0.005}{eV/\angstrom} \\
Bail-late tolerance & $10\times f_\text{max}$ \\
\bottomrule
\end{tabular}

Energy-based direction assignment is non-trivial: without it, the same reaction discovered from $A$ vs.\ $B$ of an $A \rightleftharpoons B$ pair would enter the graph in opposite orientations, and downstream dedup would have to handle the symmetric case explicitly. The bail-late tolerance is generous relative to the convergence criterion but appropriate at ML level; the DFT refinement (\S\ref{sec:kinetics-dft}) supersedes the ML barriers consumed by the kinetics solver downstream.


%-----------------------------------------------------------------------------
\subsubsection{Stage 6 --- Endpoint identity}
\label{sec:endpoints}
%-----------------------------------------------------------------------------

\paragraph{Strict mode.}
Exactly one endpoint must match the known reactant, with the match defined by permutation-invariant RMSD (\S\ref{sec:shared-rmsd}) below $\SI{1.0}{\angstrom}$. For single-molecule contexts the reactant is the source compound; for merge contexts the endpoint is fragmented (\S\ref{sec:shared-fragments}) and the validator tries both fragment-to-component assignments. Three outcomes: \textsc{one match} (accept), \textsc{both match} (no new product), \textsc{neither match} (fall through to greedy).

\paragraph{Greedy fallback.}
Enabled by default for both merge and single contexts. The validator drops the source-compound requirement and accepts the reaction iff the two endpoints describe chemically different states: each side fragmented into $\leq 2$ components with canonically distinct SMILES (\S\ref{sec:shared-smiles}) multisets, and no \emph{spectator species} (a fragment that appears on both sides at identical multiplicity while the other species differs --- the reaction $A + B \to A + C$ is rejected because $A$ is unchanged, but $2A \to A + B$ is accepted because $\{A,A\} \neq \{A,B\}$).

Without greedy mode, $> 80\%$ of valid TS proposals would be rejected because the IRC found a reaction whose reactant differs from the sampler's pick (different conformer, tautomer, or related-but-distinct species). Greedy-accepted reactions register against the actual IRC-determined reactants, which may differ from the seeding context's compounds.

\noindent\begin{tabular}{lr}
\toprule
Parameter & Value \\
\midrule
RMSD match threshold (strict) & \SI{1.0}{\angstrom} \\
Greedy fallback (single + merge) & enabled \\
Max fragments / side (greedy) & 2 \\
\bottomrule
\end{tabular}

The \SI{1.0}{\angstrom} threshold is much looser than the \SI{0.20}{\angstrom} minimum-deduplication cut used at the graph-commit stage because the test here is ``same chemical species, possibly different conformer'' --- conformer-collapse in the per-compound PES graph handles the disambiguation downstream. Greedy mode is the single most consequential pipeline choice: it allows the diffusion model to discover reactions the sampler did not anticipate, but it imports any biases of the model.


%-----------------------------------------------------------------------------
\subsubsection{Stage 7 --- Fragmentation detection}
\label{sec:fragmentation}
%-----------------------------------------------------------------------------

The product side is decomposed into connected components by the covalent-radius-scaled distance criterion (\S\ref{sec:shared-fragments}), with the standard validity rules (non-empty; satisfies the tetrose composition cap; single-atom only if hydrogen). The fragmentation is all-or-nothing: if any component fails validity, the context is dropped --- partial validity would yield inconsistent reaction stoichiometry.

For each valid fragment, the canonical SMILES is computed and an integer charge is inferred from the perceived bond pattern (\S\ref{sec:shared-smiles}); the DFT refinement (\S\ref{sec:kinetics-dft}) re-derives charges from a wavefunction calculation downstream. When a product fragments into multiple components, each fragment is independently re-relaxed on its own PES (FIRE, \S\ref{sec:shared-fire}; $f_\text{max} = \SI{0.005}{eV/\angstrom}$, max $200$ steps) to remove the residual inter-fragment distortion left by the combined-system IRC.

\noindent\begin{tabular}{lr}
\toprule
Parameter & Value \\
\midrule
Covalent-radius scale & 1.3 \\
Per-fragment composition cap & $\{C{:}5,H{:}10,O{:}5\}$ \\
Single-atom rule & H only \\
Per-fragment $f_\text{max}$ & \SI{0.005}{eV/\angstrom} \\
Per-fragment max steps & 200 \\
\bottomrule
\end{tabular}

Per-fragment relaxation matters because the IRC product is a minimum of the combined (interacting) system, not of each fragment in isolation; un-corrected, the per-fragment energies would not be comparable to those of the same fragments encountered in other reactions.


%-----------------------------------------------------------------------------
\subsubsection{Stage 8 --- Backward barrier gate}
\label{sec:bwd-barrier}
%-----------------------------------------------------------------------------

Symmetric to Stage 4. The backward barrier is $\Delta E_\text{bwd} = E_\text{TS} - E_\text{product}$ (summed over fragments if applicable), accepted iff $\Delta E_\text{bwd} \in (\SI{0}{eV}, \SI{2.7211}{eV}]$, with \emph{strict} positivity (a negative backward barrier means a non-stationary or non-converged IRC). The barrier is computed by trajectory force integration (\S\ref{sec:shared-traj}). Surviving contexts pass to the graph-addition step, which canonicalises and registers reactant/product compounds, deduplicates minima against the per-compound PES graphs (three-tier RMSD + energy fast-path, then NEB at \SI{0.05}{eV} barrier threshold; \S\ref{sec:shared-rmsd}, \S\ref{sec:shared-neb}), and either creates a new reaction row or updates an existing lower-barrier variant. Distinct sampling contexts that converged to the same reaction collapse at this point.


%-----------------------------------------------------------------------------
\subsubsection{Pipeline funnel statistics}
\label{sec:funnel}
%-----------------------------------------------------------------------------

\begin{table}
\centering
\begin{tabular}{lrr}
\toprule
Stage & Survivors & Cumulative survival \\
\midrule
Submitted (Stage 1)                                  & $3\,814\,199$         & $100.00$\% \\
After Stages 2--3 (merge prep + diffusion + dedup)   & $3\,814\,199$         & $100.00$\% \\
After Stage 4 (fwd barrier)                          & $1\,798\,768$         & $47.16$\%  \\
After Stages 5--6 (Hessian + IRC + endpoint id)      & $172\,267$            & $4.52$\%   \\
After Stage 7 (fragmentation)                        & $172\,267$            & $4.52$\%   \\
After Stage 8 (bwd barrier)                          & $127\,670$            & $3.35$\%   \\
\midrule
Added to graph (post-dedup)                          & $30\,493$             & $0.80$\%   \\
\bottomrule
\end{tabular}
\caption{Generative-loop funnel for the main run, aggregated over all batches. For comparison, the PES loop (SI-\ref{si:pes}) added $16\,879$ of $69\,450$ escaped reactions, so the generative loop contributed $\sim 64\%$ of the reactions in the graph.}
\label{tab:funnel}
\end{table}

\paragraph{Attrition.}
\textbf{Stage 4} kills $\sim 53\%$ of denoised candidates; almost all rejections are above the upper bound ($> \SI{2}{eV}$ overflow: $1{,}093{,}685$ merge + $1{,}050{,}310$ single), corresponding to strained or partially dissociated denoiser outputs. \textbf{Stages 5--6} kill $> 90\%$ of the remainder: of $1{,}798{,}768$ entering IRC, $418{,}475$ fail the Hessian check; $125{,}090$ and $89{,}898$ fail forward/backward relaxation; $706{,}221$ have both endpoints relax back to the source (\textsc{both\_match}); $412{,}454$ have neither match; only $46{,}630$ pass strict matching. Greedy fallback engages on the $1{,}118{,}675$ strict-mismatched cases and accepts $125{,}637$ (the dominant greedy rejection is identical SMILES multisets on both sides: $961{,}323$). \textbf{Stage 7} rejected zero contexts in the main run --- every IRC-passing geometry fragmented cleanly. \textbf{Stage 8} kills a further $44{,}597$ ($\sim 26\%$) on near-zero backward barriers. \textbf{Graph dedup} removes $97{,}177$ ($\sim 76\%$ of post-Stage-8 survivors), reflecting the fact that many distinct sampling contexts converge to the same reaction.

\paragraph{Merge vs.\ single.}
Submitted: $2{,}613{,}816$ single, $1{,}200{,}383$ merge. After Stage 4: $1{,}706{,}049$ single ($65.3\%$) vs.\ $92{,}719$ merge ($7.7\%$) --- merged complexes are far more often rejected at the cheap energy gate because the fragment-sum baseline is much further from the TS. After Stage 5: $164{,}070$ vs.\ $8{,}197$ ($\sim 9\times$ asymmetric). Greedy-accepted: $120{,}808$ single vs.\ $4{,}829$ merge --- single contexts depend on greedy almost universally, merges mostly succeed under strict matching.

% --- end si2_generative_loop.tex ---

% --- begin si1_pes_loop.tex ---
% ============================================================================
% Supplementary Information SI-1: PES-Loop Reaction Discovery
%
% Drop-in subsection fragment. Parent SI must provide: amsmath/amssymb,
% booktabs, siunitx (with \angstrom). Bibliography handled at the parent
% level. (Parameter reference lives in si_parameters.tex.)
% ============================================================================

\subsection{PES-Loop Reaction Discovery}
\label{si:pes}

The PES loop performs local exploration around a known minimum. It works one compound at a time, MD-samples configuration space near a chosen minimum, optimises trajectory frames to first-order saddle points by Partitioned Rational Function Optimisation (P-RFO), and validates each saddle by IRC-like descent. The endpoints of that descent fall into two cases: \emph{intramolecular} (both endpoints belong to the same compound's PES graph --- a conformer-to-conformer reaction) or \emph{escape} (one endpoint is a structurally different compound, possibly fragmented into multiple species --- a true chemical reaction). Intramolecular saddles are stored as edges in the per-compound PES graph; escapes are passed to the same multi-stage critic used by the generative loop (SI-\ref{si:generative}, Stages 4--8) and admitted to the reaction graph if they survive.

The pipeline has eight stages: (1)~minimum selection, (2)~molecular dynamics, (3)~TS candidate sampling, (4)~batched P-RFO, (5)~Hessian saddle validation, (6)~bidirectional IRC, (7)~endpoint classification (intramolecular vs.\ escape), (8)~commitment. Energies, forces, and Hessians come from the ML force field; that model and the shared numerical primitives (Hessian projection, FIRE, trust-region Newton, IRC, NEB, RMSD, fragmentation, SMILES) are documented once in SI-\ref{si:shared}.


%-----------------------------------------------------------------------------
\subsubsection{Stage 1 --- Minimum selection}
\label{sec:pes-select}
%-----------------------------------------------------------------------------

PES exploration is queued at the granularity of \emph{(compound, minimum)} pairs. A worker claims one minimum at a time via a per-compound advisory lock (so only one worker explores a given compound at any moment). Within a claimed compound, the lowest-energy unexplored minimum is selected:
\begin{equation}
    m_\text{next} = \arg\min_{m \in \mathcal{M}_\text{unexplored}(c)} E(m).
\end{equation}
Lower-energy conformers are more populated and host the kinetically dominant reactions, so they are explored first. Once selected, the minimum is marked in-progress; on success it is marked \emph{explored} and the worker releases the compound. Compounds whose every minimum has been explored leave the queue.


%-----------------------------------------------------------------------------
\subsubsection{Stage 2 --- Molecular dynamics}
\label{sec:pes-md}
%-----------------------------------------------------------------------------

A Langevin trajectory is launched from the selected minimum to sample configuration space around it. The goal is not kinetics simulation but geometric coverage --- to visit a diverse set of geometries that lie on or near reaction pathways departing from the minimum. The integrator follows
\begin{equation}
    m_i \ddot{\mathbf{r}}_i = \mathbf{F}_i - \gamma m_i \dot{\mathbf{r}}_i + \sqrt{2 \gamma m_i k_B T}\, \boldsymbol{\xi}(t),
\end{equation}
with friction coefficient $\gamma = 1/\tau$, coupling time $\tau = \SI{100}{fs}$, and Gaussian white noise $\boldsymbol{\xi}$. Initial velocities are drawn from the Maxwell--Boltzmann distribution at the target temperature. $2500$ steps $\times \SI{0.5}{fs} = \SI{1.25}{ps}$ at \SI{300}{K}, with a snapshot saved every $10$ steps ($250$ frames).

\noindent\begin{tabular}{lr}
\toprule
Parameter & Value \\
\midrule
MD steps & 2500 \\
Integration step & \SI{0.5}{fs} \\
Total simulation time & \SI{1.25}{ps} \\
Temperature & \SI{300}{K} \\
Thermostat & Langevin \\
Coupling $\tau$ & \SI{100}{fs} \\
Snapshot interval & 10 steps (250 frames) \\
\bottomrule
\end{tabular}

At \SI{300}{K}, $k_B T \approx \SI{0.026}{eV}$ --- direct sampling of saddle points with barriers above $\sim \SI{0.5}{eV}$ is statistically negligible in \SI{1.25}{ps}. The MD is explicitly not asked to cross barriers; P-RFO does that downstream. The \SI{0.5}{fs} timestep is safe for the fastest vibrations (C--H stretches, period $\sim \SI{10}{fs}$); the moderately tight \SI{100}{fs} coupling prevents kinetic-energy build-up in any single mode without aggressively damping motion.


%-----------------------------------------------------------------------------
\subsubsection{Stage 3 --- TS candidate sampling}
\label{sec:pes-candidates}
%-----------------------------------------------------------------------------

From the 250 MD frames, $32$ candidates are selected by uniform random sampling. P-RFO is robust to starting-point choice within the same basin --- 32 random frames cover the relevant saddles on the boundary of a low-energy basin, and subsequent deduplication absorbs redundant starting points that converge to the same saddle.


%-----------------------------------------------------------------------------
\subsubsection{Stage 4 --- Partitioned Rational Function Optimisation (P-RFO)}
\label{sec:pes-prfo}
%-----------------------------------------------------------------------------

P-RFO (Banerjee et al.\ 1985; Baker 1986) walks each candidate geometry uphill along one mode and downhill in the orthogonal subspace, converging to a first-order saddle. The pipeline runs all 32 candidates in lockstep (one batched GPU Hessian + one batched energy evaluation per step via \S\ref{sec:shared-mdet}). At each step:

\begin{enumerate}
    \item Compute $E$, $\mathbf{f} = -\nabla E$, and the projected Hessian $\mathbf{H}_\text{proj}$ (eigenpairs $\{\lambda_i, \mathbf{u}_i\}$; \S\ref{sec:shared-hessian}). Project the gradient: $g_i = \mathbf{u}_i \cdot \mathbf{g}$.
    \item Compute the RFO step in each mode,
    \begin{equation}
        p_i = \frac{\mu_i}{g_i}, \qquad \mu_i = \frac{\lambda_i \pm \sqrt{\lambda_i^2 + 4 g_i^2}}{2},
    \end{equation}
    with $+$ on the TS mode (walk uphill) and $-$ on the others (walk downhill).
\end{enumerate}

\paragraph{Mode following.}
At step $0$ the TS mode is the lowest-eigenvalue mode. At subsequent steps it is the eigenvector with the largest overlap to the previous TS eigenvector,
\begin{equation}
    i_\text{TS} = \arg\max_i \left| \mathbf{u}_i \cdot \mathbf{u}_\text{TS}^{(\text{prev})} \right|.
\end{equation}
This prevents the optimiser from jumping between modes when eigenvalues cross.

\paragraph{Partitioned trust radii.}
Two adaptive trust radii are maintained: a TS-mode radius (initially $1.5\times$ the minimum-mode radius) and a minimisation-subspace radius. After each trial step, the trust-region ratio
\begin{equation}
    \rho = \frac{\Delta E_\text{actual}}{\Delta E_\text{predicted}}, \qquad \Delta E_\text{pred} = -\mathbf{f} \cdot \mathbf{p} + \tfrac{1}{2} \mathbf{p}^\top \mathbf{H} \mathbf{p}
\end{equation}
is computed. The step is accepted if $\rho > 0.1$ (otherwise the trust radius is halved and the step retried, up to $5$ times). After acceptance, the radius shrinks by $0.5$ if $\rho < 0.25$ or $\rho > 2$ (poor model fit), and expands by $1.5$ if $0.5 < \rho < 1.5$ and the step hit the trust boundary.

\paragraph{Convergence and early stopping.}
A candidate is declared converged when $\max_i |f_i| < \SI{0.01}{eV/\angstrom}$ \emph{and} $f_\text{RMS} < \SI{0.005}{eV/\angstrom}$ \emph{and} the projected Hessian has exactly one significant negative eigenvalue. The optimiser also detects pathological trajectories: \emph{stuck} (net displacement and total path length below $10^{-4}$~\AA{} over the last 20 steps), \emph{oscillating} (net/path ratio $< 0.1$), or \emph{bad steps} ($>20$ consecutive steps with $\rho < 0.1$ or $\rho > 2$ at minimum trust radius).

\noindent\begin{tabular}{lr}
\toprule
Parameter & Value \\
\midrule
Candidates per iteration & 32 \\
Max steps & 120 \\
$f_\text{max}$ tolerance & \SI{0.01}{eV/\angstrom} \\
$f_\text{RMS}$ tolerance & \SI{0.005}{eV/\angstrom} \\
Initial trust radius & \SI{0.5}{\angstrom} \\
Trust radius bounds & $[0.01, 0.3]$~\AA \\
TS-mode trust scale & 1.5 \\
Accept threshold ($\rho$) & 0.1 \\
Max retries / step & 5 \\
\bottomrule
\end{tabular}

The \SI{0.01}{eV/\angstrom} convergence is tight enough to be meaningful at ML force-field accuracy but loose enough to actually converge --- tighter would just chase model noise. The permissive $\rho > 0.1$ accept threshold reflects that the ML Hessian sometimes misses small higher-order terms, and a step that disagrees with the quadratic model by up to 10$\times$ may still be productive.


%-----------------------------------------------------------------------------
\subsubsection{Stage 5 --- Hessian saddle validation}
\label{sec:pes-saddle-validation}
%-----------------------------------------------------------------------------

A P-RFO-converged geometry may have spurious extra negative eigenvalues from the ML model (especially in floppy directions like methyl rotations or ring puckers) or, conversely, a single negative eigenvalue too shallow to support a real reaction. The projected Hessian (\S\ref{sec:shared-hessian}) is recomputed at the converged geometry and the spectrum examined. A saddle is admitted iff:
\begin{itemize}
    \item \emph{exactly one} eigenvalue satisfies $\lambda < -10^{-4}\,\text{eV/\AA}^2$ (significance threshold, masking numerical noise), and
    \item the most negative eigenvalue satisfies $\lambda_\text{TS} < \SI{-0.01}{eV/\angstrom^2}$ (the curvature is meaningfully negative, not just below numerical floor).
\end{itemize}
Candidates with $-0.01 < \lambda_\text{TS} < 0$ are rejected as ``eigenvalue too shallow''. The two-tier check separates two failure modes: noise pretending to be a saddle (lowest eigenvalue just below zero from numerical floor) and insufficient curvature to call a saddle (a real but very shallow negative eigenvalue). The exactly-one criterion is strict by design --- the PES loop submits only candidates its P-RFO walker drove to a clean first-order saddle, so allowing additional negative modes here would admit converged-but-pathological geometries that should be discarded rather than salvaged. (The generative loop deliberately relaxes this to ``at least one'' in \S\ref{sec:hessian-irc}, because diffusion-proposed geometries span a much broader curvature distribution and downstream IRC + barrier gates absorb the additional load.) ML Hessians for very floppy systems can fail this check on geometries that DFT would confirm as saddles; for the downstream kinetics that ultimately consume the barriers, this is mitigated by the DFT energy refinement (\S\ref{sec:kinetics-dft}).


%-----------------------------------------------------------------------------
\subsubsection{Stage 6 --- Bidirectional IRC}
\label{sec:pes-irc}
%-----------------------------------------------------------------------------

The imaginary mode is followed in both directions and relaxed to a minimum on each side; the bidirectional IRC algorithm is shared with the generative loop and described in \S\ref{sec:shared-irc}. PES-loop-specific choices:

\paragraph{Curvature-adaptive displacement.}
The harmonic-approximation form targets an energy drop of $\Delta E_\text{target} = 2 \times \SI{0.04}{eV}$ (twice the endpoint validation threshold below):
\begin{equation}
    \Delta x = \sqrt{\frac{2|\Delta E_\text{target}|}{|\lambda_\text{TS}|}}, \qquad \text{clamped to } [\SI{0.02}{\angstrom}, \SI{0.5}{\angstrom}].
\end{equation}
Endpoint relaxation uses trust-region Newton (\S\ref{sec:shared-newton}) with $f_\text{max} = \SI{0.005}{eV/\angstrom}$, max $120$ steps, and an energy-increase tolerance of $\SI{0.005}{eV}$ per step. Acceptance: convergence or bail-late at $f_\text{max} < 10\times$ tolerance.

\paragraph{Endpoint criteria.}
Both endpoints must lie energetically below the TS. The standard criterion is symmetric:
\begin{equation}
    E_\text{left} - E_\text{TS} < \SI{-0.04}{eV} \quad \text{and} \quad E_\text{right} - E_\text{TS} < \SI{-0.04}{eV}.
\end{equation}
An asymmetric relaxation is also accepted, in which one side passes the standard threshold and the other a relaxed threshold:
\begin{equation}
    \bigl(E_\text{left} < E_\text{TS} - 0.04 \text{ and } E_\text{right} < E_\text{TS} - 0.005\bigr) \quad \text{or vice versa},
\end{equation}
which catches genuine early/late TSs whose product- or reactant-side well is very shallow. Both endpoints must be structurally distinct, with permutation-invariant RMSD $> \SI{0.25}{\angstrom}$ (otherwise the IRC found two near-identical structures and the ``saddle'' was a numerical artefact). The higher-energy endpoint is labelled \emph{forward} (reactant), the lower-energy \emph{backward} (product), so the reaction is consistently stored downhill in the PES-graph view.

\noindent\begin{tabular}{lr}
\toprule
Parameter & Value \\
\midrule
Energy-drop target & \SI{0.08}{eV} \\
Displacement formula & $\sqrt{2 \cdot 0.08 / |\lambda|}$~\AA \\
Displacement clamp & $[0.02, 0.5]$~\AA \\
Endpoint $f_\text{max}$ & \SI{0.005}{eV/\angstrom} \\
Endpoint max steps & 120 \\
Endpoint energy threshold & \SI{0.04}{eV} below TS (both sides) \\
Asymmetric threshold & \SI{0.005}{eV} (one side only) \\
Min endpoint RMSD & \SI{0.25}{\angstrom} \\
\bottomrule
\end{tabular}

The asymmetric \SI{0.005}{eV} relaxation is necessary for late/early TSs where one side of the saddle is very flat --- without it, ring-opening and similar very-asymmetric reactions would be silently lost. However, $\SI{0.005}{eV} \approx \SI{0.12}{kcal/mol}$ is close to ML-model noise on energies; reactions accepted only by the asymmetric criterion deserve DFT cross-validation. The \SI{0.25}{\angstrom} RMSD floor catches IRCs that collapsed back into the same conformer but does not enforce ``chemically distinct'' --- that is checked at Stage 7 below and again at the per-compound minimum dedup (\S\ref{sec:shared-rmsd}, \S\ref{sec:shared-neb}).


%-----------------------------------------------------------------------------
\subsubsection{Stage 7 --- Endpoint classification: intramolecular vs.\ escape}
\label{sec:pes-classify}
%-----------------------------------------------------------------------------

After IRC, each endpoint is canonicalised and a SMILES is inferred from its bonded structure (\S\ref{sec:shared-smiles}). Three classifications follow.

\paragraph{Intramolecular.}
Both endpoints have SMILES matching the originating compound. The TS connects two minima of the same compound's PES (different conformers, different protomers within the same connectivity, or different ring puckers). The saddle is committed as an edge in the per-compound PES graph (Stage 8a).

\paragraph{Escape.}
Exactly one endpoint matches the originating compound's SMILES; the other does not. The IRC has discovered a reaction that leaves the originating compound's basin and lands in a chemically different species (possibly multiple, if the relaxed product fragmented). The reaction is repackaged as an \texttt{ExplorationContext} with the originating compound as reactant and the escaped structure as product, and passed to the same multi-stage critic as the generative loop (SI-\ref{si:generative}, Stages 4--8): forward barrier gate, Hessian + IRC validation (re-run on the relocated TS for consistency), endpoint identity (greedy mode, since the escape product is by construction not in the source compound's PES), fragmentation detection, backward barrier gate.

\paragraph{Anomalous.}
Neither endpoint matches the originating compound's SMILES, or SMILES inference fails. The saddle is discarded. These are uncommon and typically reflect IRC drift into chemistry that the relaxation cannot ground: a transient charge-separated state, a partially formed bond that RDKit cannot perceive, or a fragmentation where one fragment failed canonicalisation.

The escape branch is what makes the PES loop a \emph{reaction-discovery} engine rather than a conformer enumerator. Most reactions accessible from a given minimum are conformational (intramolecular) and stay within one compound's PES graph; a smaller fraction cross to new chemistry. In the main run, the PES loop produced $23\,720$ intramolecular saddles versus $69\,450$ escape candidates, of which $16\,879$ ($\sim 24\%$) survived the generative critic (SI-\ref{si:generative}, Stages 4--8).


%-----------------------------------------------------------------------------
\subsubsection{Stage 8 --- Commitment}
\label{sec:pes-commit}
%-----------------------------------------------------------------------------

\paragraph{8a. Intramolecular saddles $\to$ per-compound PES graph.}
The saddle and its two endpoints are added to the per-compound graph (one PES graph per compound). New endpoint minima go through the three-tier minimum deduplication: RMSD (\S\ref{sec:shared-rmsd}) $<$ \SI{0.01}{\angstrom} ultra-low merge; RMSD $< \SI{0.20}{\angstrom}$ and $|\Delta E| < \SI{0.05}{eV}$ fast merge; otherwise NEB same-basin check (\S\ref{sec:shared-neb}) at \SI{0.05}{eV} barrier threshold. The saddle is added as an edge unless an existing edge between the same endpoint pair has matching barrier heights within \SI{0.1}{eV} (TS edge dedup). Edge barriers prefer the trajectory force-integral form (\S\ref{sec:shared-traj}) over the endpoint difference.

\paragraph{8b. Escapes $\to$ reaction graph.}
Surviving escape contexts (after SI-\ref{si:generative} Stages 4--8) are added to the reaction graph with \texttt{discovery\_method = "pes\_exploration"}, distinguishing them from the generative-loop reactions tagged \texttt{"generative"} and the manual buffer equilibria tagged \texttt{"manual\_equilibrium"} (SI-\ref{si:seed}). Compound registration, reaction-multiset deduplication, and barrier-source priority (DFT $>$ ML in-box; SI-\ref{si:kinetics}) are identical to the generative loop.


%-----------------------------------------------------------------------------
\subsubsection{Pipeline funnel statistics}
\label{sec:pes-funnel}
%-----------------------------------------------------------------------------

\begin{table}
\centering
\begin{tabular}{lr}
\toprule
Quantity & Count \\
\midrule
PES jobs run (minima explored)                        & $26\,282$ \\
Intramolecular saddles found and committed            & $23\,720$ \\
Escape candidates produced                            & $69\,450$ \\
Escape candidates surviving generative critic (SI-\ref{si:generative}, Stages 4--8) & $16\,879$ \\
Escape critic survival rate                           & $\sim 24\%$ \\
\bottomrule
\end{tabular}
\caption{PES-loop totals for the main run. The escape critic survival rate ($24\%$) is much higher than the generative-loop submission-to-graph rate ($0.8\%$, SI-\ref{si:generative}) because the PES-loop ``proposals'' are P-RFO-converged geometries with a validated Hessian spectrum and a successful bidirectional IRC --- i.e., they have already passed the equivalent of generative Stages 5--6 before entering the critic.}
\label{tab:pes-funnel}
\end{table}

The PES loop contributed $\sim 36\%$ of all reactions in the graph ($16\,879$ of the $\sim 47\,000$ combined). It dominates two regions of chemical space where the generative loop is weak: (i)~conformational transitions within a compound (the entire intramolecular-saddle population, which the generative loop cannot produce because it does not sample within-compound saddles), and (ii)~reactions whose TS geometry is far from the diffusion model's training distribution but is reachable by MD from a known minimum. The PES loop's reach is bounded by what MD samples at \SI{300}{K} in \SI{1.25}{ps} --- it cannot find reactions whose reactants are not already in the compound table.

% --- end si1_pes_loop.tex ---

% --- begin si3_kinetics.tex ---
% ============================================================================
% Supplementary Information SI-3: Reaction-Network Kinetics Solver
%
% Drop-in subsection fragment. Parent SI must provide: amsmath/amssymb,
% booktabs, siunitx (with \angstrom), and the tcolorbox environments
% {evaluation}, {concern}, {threshbox}. Bibliography handled at the
% parent level. (Parameter reference lives in si_parameters.tex.)
% ============================================================================

\subsection{Reaction-Network Kinetics Solver}
\label{si:kinetics}

The kinetics solver consumes the current reaction graph and integrates a mass-action ODE to a long-time steady state. The resulting steady-state composition is the distribution used by the generative-loop sampler (SI-\ref{si:generative}).

%-----------------------------------------------------------------------------
\subsubsection{Mass-action ODE}
\label{sec:kinetics-ode}
%-----------------------------------------------------------------------------

For $n_\text{species}$ species and $n_\text{reactions}$ reactions, the ODE is the standard mass-action system
\begin{equation}
    \frac{d\mathbf{y}}{dt} = \mathbf{S}_\text{net}\,\mathbf{r}(\mathbf{y}), \qquad r_j(\mathbf{y}) = k_j^+ \prod_i y_i^{\nu_{ij}^-} - k_j^- \prod_i y_i^{\nu_{ij}^+},
\end{equation}
where $\mathbf{S}_\text{net} = \mathbf{S}^+ - \mathbf{S}^-$ is the net stoichiometry matrix and $\boldsymbol{\nu}^\pm$ are the reactant/product stoichiometry exponents. The right-hand side and its analytical sparse Jacobian
\begin{equation}
    J_{ik} = \frac{\partial}{\partial y_k}\!\left[\sum_j S_{ij}\, r_j(\mathbf{y})\right]
\end{equation}
are implemented as Numba-JIT kernels (\texttt{numba\_kernels.py}, \texttt{@njit(cache=True)}) to keep the inner solver loop in compiled code. Jacobian sparsity follows from the stoichiometry pattern: $J_{ik} \neq 0$ only where species $k$ appears as a reactant or product of a reaction that touches species $i$.

%-----------------------------------------------------------------------------
\subsubsection{Rate constants}
\label{sec:kinetics-rates}
%-----------------------------------------------------------------------------

\paragraph{Eyring transition-state theory.}
For each discovered reaction the forward and backward rate constants are
\begin{equation}
    k = \min\!\Bigl(\tfrac{k_B T}{h}\,\exp(-E_a / k_B T),\ k_\text{diff}\Bigr), \qquad k_\text{diff} = \SI{e10}{M^{-1}s^{-1}},
\end{equation}
with $k_B = \SI{8.617e-5}{eV/K}$, $h = \SI{4.136e-15}{eV\,s}$, and the input barrier capped at $E_a \leq \SI{10}{eV}$ before the exponential (purely to avoid floating-point overflow; barriers above this are off-policy anyway). The diffusion-limit cap prevents the unit-prefactor Eyring rate from exceeding the physical bimolecular collision rate at the default exploration temperature of \SI{500}{K} --- without it, near-zero barriers would produce non-physical rates that dominate the Jacobian conditioning.

\paragraph{Barrier source priority.}
The barrier that feeds Eyring is chosen by a strict priority:
\begin{enumerate}
    \item PBE0/def2-TZVPP separated barriers (\texttt{barrier\_\{forward,backward\}\_separated\_pbe0}) if both forward and backward are available --- this is the DFT-refined ground truth from the CPU pipeline (\S\ref{sec:kinetics-dft}).
    \item Force-integrated in-box ML barriers (\texttt{barrier\_\{forward,backward\}}) along the IRC trajectory.
    \item Otherwise: the reaction is dropped entirely from the ODE.
\end{enumerate}
\emph{ML separated-energy barriers (computed as the difference of energy-head evaluations on two unrelated geometries) are never used} --- they are the least reliable variant: differencing two energy-head evaluations on unrelated geometries is noisier than either the DFT separated estimate or the ML force-integrated estimate.

\paragraph{Manual buffer equilibria.}
Acid--base equilibria (water autoionisation, CO$_2$ hydration, H$_2$CO$_3$ dissociation, proton solvation; SI-\ref{si:seed}) bypass Eyring entirely. Their rate constants are stored as literal values in the \texttt{Reaction} row (\texttt{manual\_k\_fwd}, \texttt{manual\_k\_bwd}) and applied directly at any temperature. These reactions are tagged \texttt{discovery\_method = "manual\_equilibrium"} and are layered into the ODE on top of the discovered reactions.

\paragraph{Filters applied at build time.}
Reactions are dropped from the ODE if (i)~either side contains dot-disconnected SMILES (van der Waals dimers stabilised by dispersion; non-chemical), (ii)~either barrier is below $-\SI{0.05}{eV}$ (negative-$E_a$ artefacts from PES/IRC landing on dianion or zwitterion minima below the separated reactant pair), or (iii)~both barriers exceed the cutoff (default \SI{10}{eV} --- effectively unreachable, just bloats the system). Reactions are also direction-agnostic deduplicated: $A + B \to C$ and $C \to A + B$ collapse into one bidirectional reaction keyed on the unordered pair of reactant/product multisets, taking the lowest $\min(E_a^\text{fwd}, E_a^\text{bwd})$. Storing both as separate mass-action reactions would double-count the net flux for the same elementary step.

The Eyring prefactor at $T = \SI{500}{K}$ gives $k_B T / h \approx \SI{1.04e13}{s^{-1}}$, so a $\SI{1.0}{eV}$ barrier corresponds to $\sim \SI{9e2}{s^{-1}}$ and a $\SI{2.0}{eV}$ barrier to $\sim \SI{8e-8}{s^{-1}}$. The \SI{500}{K} default is a sampling choice: at the higher temperature, more reactions are kinetically accessible within the integration window, producing a richer steady-state distribution for the generative sampler to draw from. Physical-temperature analyses use the same graph but a separate solve at the target temperature.


%-----------------------------------------------------------------------------
\subsubsection{DFT energy refinement}
\label{sec:kinetics-dft}
%-----------------------------------------------------------------------------

The sole purpose of the DFT pipeline in this system is to upgrade the barriers fed to the Eyring rates above. There is no DFT geometry optimisation, no DFT vibrational analysis, and no DFT-based dataset assembly: only single-point energies on the geometries already produced by the ML pipeline, used to derive PBE0 barriers that replace the ML in-box estimates whenever available.

\paragraph{Method.}
PBE0/def2-TZVPP single-points via PySCF (\texttt{dft.RKS} for closed shell, \texttt{dft.UKS} for open shell; spin guessed from electron count modulo 2). SCF convergence \texttt{conv\_tol} $= 10^{-9}$, \texttt{max\_cycle} $= 300$, with a Newton retry on non-convergence. All energies are returned in eV.

\paragraph{Per-reaction work.}
For each non-manual reaction the pipeline computes three single-point energies on the in-box geometry (same atom set throughout, since the IRC is performed in the merged-system box):
\begin{itemize}
    \item $E_\text{TS}^\text{DFT}$ at the TS conformer geometry,
    \item $E_R^\text{DFT}$ at the relaxed reactant frame (the final position of the IRC reactant-side trajectory),
    \item $E_P^\text{DFT}$ at the relaxed product frame (final position of the IRC product-side trajectory).
\end{itemize}
The in-box DFT barriers are then $\Delta E_\text{fwd}^\text{in-box} = E_\text{TS} - E_R$ and $\Delta E_\text{bwd}^\text{in-box} = E_\text{TS} - E_P$.

\paragraph{Per-compound reference energies.}
For each compound participating in any DFT-refined reaction, a single-point at the compound's \emph{lowest-energy minimum} (the PES-graph entry with the smallest ML energy) is computed and cached on the \texttt{Compound} row. The cache is keyed by \texttt{(method, basis)}; subsequent reactions touching the same compound reuse the value rather than recomputing. The \emph{separated} barriers --- the ones the kinetics solver actually consumes --- are derived as
\begin{equation}
    \Delta E_\text{fwd}^\text{sep} = E_\text{TS}^\text{DFT} - \!\!\sum_{c \in \text{reactants}}\!\! E_c^\text{DFT}, \qquad
    \Delta E_\text{bwd}^\text{sep} = E_\text{TS}^\text{DFT} - \!\!\sum_{c \in \text{products}}\!\! E_c^\text{DFT},
\end{equation}
where the sums respect reactant/product stoichiometry. Because the reference energies use the compound's lowest minimum (not the IRC-relaxed in-box geometry), the separated barriers absorb the inter-fragment binding-energy correction that the in-box variant misses.

\paragraph{No Hessians, no frequency check.}
The DFT pipeline does not compute Hessians or check TS character at DFT level. The reasoning is empirical: hybrid-functional Hessians fail on the same floppy-mode pathologies as ML Hessians (small spurious negative eigenvalues from low-frequency modes, near-degenerate hindered rotors), at $10^2$--$10^3\times$ the cost. Energies are where DFT provides leverage over ML; for spectrum verification, a dedicated tighter ML model or a CC-level audit on a small subset would be more cost-effective than per-reaction DFT Hessians.

\paragraph{Cost.}
A PBE0/def2-TZVPP single-point scales as $\mathcal{O}(N^{2}$--$N^{3})$ with system size; wall times in the production deployment range from $\sim$\SI{10}{s} for diatomic fragments to $\sim$\SI{5}{min} per geometry for $N \sim 13$--$20$ heavy atoms with the def2-TZVPP basis. Each reaction therefore costs three TS-sized single-points plus (amortised over the compound population) a handful of reference single-points. The pipeline runs on a separate pool of CPU workers (one work item $=$ one reaction) that drain a Postgres-backed queue independently of the GPU exploration workers, so DFT throughput does not bottleneck reaction discovery.

Refining only energies, on only the geometries the ML model already produced, is a deliberate cost--benefit trade. Re-optimising at DFT level would correct ML geometry errors but add $10$--$50\times$ the wall time per reaction and would require a DFT-level IRC re-run to keep the reactant/product matchings consistent. Empirically the ML geometry errors are small for the chemistry covered here (RMSD $\lesssim \SI{0.05}{\angstrom}$ to PBE0 minima for relaxed carbohydrate-space compounds), so the single-point-only protocol captures the bulk of the DFT correction at a small fraction of the cost. The compound-reference cache is load-bearing: without it, every reaction would re-compute energies for the same recurring small compounds (water, formaldehyde, etc.).


%-----------------------------------------------------------------------------
\subsubsection{Species universe and initial conditions}
\label{sec:kinetics-ic}
%-----------------------------------------------------------------------------

The species set is the union of (i)~every compound that appears as a reactant or product in any selected reaction, and (ii)~the seed species with prescribed initial concentrations ($\mathrm{H_2O}$, $\mathrm{CH_2O}$, $\mathrm{CO_2}$, $\mathrm{H_2}$, $\mathrm{H^+}$, $\mathrm{OH^-}$, $\mathrm{H_2CO_3}$, $\mathrm{HCO_3^-}$, $\mathrm{H_3O^+}$). Seeds are included even if no reaction touches them, so the buffer equilibria can self-consistently relax to their target pH without depending on whether a discovered reaction happens to mention water.

Every species in the system receives a uniform baseline of $c_0 = \SI{e-3}{M}$ at $t = 0$, with the seed species overriding the baseline (all at $\SI{e-3}{M}$ in the production deployment). The uniform baseline is essential: a runtime-discovered species starting at exactly zero would never react, and the steady-state distribution would be dominated entirely by the seeds.

A species-concentration noise floor of \SI{e-20}{M} is applied throughout post-processing: anything below this is treated as ``not present'' for entropy and distribution computations.

%-----------------------------------------------------------------------------
\subsubsection{Solver}
\label{sec:kinetics-solver}
%-----------------------------------------------------------------------------

The ODE is integrated to $t_\text{max} = \SI{e8}{s}$ ($> 3$ years --- well past any realistic reactor residence time) using PETSc's BDF time-stepper with Numba-JIT RHS and Jacobian kernels.

\paragraph{PETSc configuration.}
\begin{itemize}
    \item Time stepper: \texttt{TS} type \texttt{bdf}, exact-final-time interpolation, max $5\times 10^7$ steps.
    \item Initial step: $\Delta t_0 = \SI{e-10}{s}$. The chemistry's fastest timescale at \SI{500}{K} is set by the diffusion limit, $\sim k_\text{diff}^{-1} = \SI{e-10}{s}$ for bimolecular rates, so the initial step matches.
    \item Tolerances: $\texttt{atol} = \SI{e-16}{}$, $\texttt{rtol} = \SI{e-12}{}$. Very tight because some species sit at concentrations far below the noise floor and the solver's atol must clear them.
    \item Newton inner solver: \texttt{SNES} with a \texttt{KSP} preconditioner-only LU factorisation (\texttt{type=preonly}, \texttt{pc=lu}). The Jacobian factorisation is the expensive step; MUMPS is the default direct solver when available.
\end{itemize}

\paragraph{Output grid.}
The solver reports the full trajectory at $400$ geometrically spaced time points covering $[\SI{e-12}{s}, \SI{e8}{s}]$ for the frontend concentration plot, plus the canonical \emph{decade grid} of $19$ points at $10^{-10}, 10^{-9}, \ldots, 10^{8}$~s used by the kinetics snapshot. The decade grid spans every kinetically relevant order of magnitude, from femtosecond chemistry to multi-year residence times.

% --- end si3_kinetics.tex ---

% --- begin si6_shared_machinery.tex ---
% ============================================================================
% Supplementary Information SI-6: Shared Geometry and Optimisation Machinery
%
% Drop-in subsection fragment. Parent SI must provide: amsmath/amssymb,
% booktabs, siunitx (with \angstrom), and the tcolorbox environments
% {evaluation}, {concern}, {threshbox}. Bibliography handled at the
% parent level. (Parameter reference lives in si_parameters.tex.)
% ============================================================================

\subsection{Shared Machinery: ML Force Field, Geometry, and Optimisation}
\label{si:shared}

The discovery loops (SI-\ref{si:pes}, SI-\ref{si:generative}) and the kinetics pipeline (SI-\ref{si:kinetics}) share a single on-line energy oracle and a small set of numerical primitives that are not specific to any one loop. This subsection consolidates them. Chemistry-specific thresholds (e.g., the asymmetric \SI{0.04}{}/\SI{0.005}{eV} IRC endpoint criterion in SI-\ref{si:pes}, or the \SI{1.0}{\angstrom} RMSD endpoint match in SI-\ref{si:generative}) stay in the calling SI; this section defines the model and the algorithms those thresholds gate.

%-----------------------------------------------------------------------------
\subsubsection{ML force field (MD-ET)}
\label{sec:shared-mdet}
%-----------------------------------------------------------------------------

All energies, forces, and Hessians used by the on-line pipeline come from a single transformer-based neural network potential, \textbf{MD-ET}. It is a \emph{direct force model}: forces are predicted as a per-atom 3-vector output head rather than obtained by differentiating an energy. Energy and force are independent heads, and the Hessian is obtained by automatic differentiation of the force head with respect to atomic positions (one row per force component, batched via \texttt{torch.vmap}). Radial cutoff for the neighbour list: $\SI{5.0}{\angstrom}$.

Three depths are released as the \texttt{variant} parameter --- \texttt{4l}, \texttt{5l}, \texttt{12l} --- with \texttt{12l} the production default. The model is trained on the QCML dataset (broad small-organic chemical space, DFT-level reference energies and forces); the transition-state subset is not represented separately, which is the structural reason the per-saddle Hessian still benefits from the trust-region machinery below and from the validation gates in SI-\ref{si:pes}/SI-\ref{si:generative}. The TS-proposer diffusion model (SI-\ref{si:generative} Stage 3) is trained on a separate TS-relaxed dataset and is independent of the MD-ET training distribution.

Distribution and loading: the model ships as the \texttt{md-et} pip package and is loaded via \texttt{md\_et.load\_calculator(variant=...)}; weights are pulled from the Hugging Face Hub on first use and the worker Dockerfile pre-warms the cache at build time. The returned object exposes the standard ASE \texttt{Calculator} interface plus a bespoke \texttt{get\_batched\_hessians} for the throughput-critical Hessian path. The default Hessian batch size is $16$; the production deployment uses $8$ to stay within L4 24~GB VRAM at typical pipeline sizes (the calculator's built-in halving recovery only applies within a single call, not across calls in a batch). A \texttt{filter\_forces=true} option enables a noise-suppression filter on the predicted forces and is on across the production deployment.

Routing every on-line ML call through a single direct-force calculator means the ML in-box barriers consumed by the kinetics solver (SI-\ref{si:kinetics}), P-RFO convergence, IRC endpoint stability, and the generative critic's barrier gates all depend on the same model. The DFT energy refinement (\S\ref{sec:kinetics-dft}) provides the independent ground truth that the kinetics ultimately consumes, so MD-ET acts as a fast filter whose role is to keep the candidate set tractable; the final rate constants depend on DFT whenever the CPU pipeline has caught up.


%-----------------------------------------------------------------------------
\subsubsection{Permutation-invariant RMSD}
\label{sec:shared-rmsd}
%-----------------------------------------------------------------------------

Comparing two molecular geometries requires aligning them in space and resolving the labelling ambiguity for chemically equivalent atoms. The system uses Kabsch alignment for rigid-body rotation and a two-strategy resolution for atom permutations:

\begin{enumerate}
    \item \textbf{Kabsch alignment.} Centre both structures at the origin and compute the optimal rotation via singular-value decomposition of the cross-covariance:
    \begin{equation}
        \mathbf{R}^* = \arg\min_{\mathbf{R} \in SO(3)} \sum_i \lVert \mathbf{r}_i^{(1)} - \mathbf{R}\,\mathbf{r}_i^{(2)} \rVert^2.
    \end{equation}
    The SVD-based form handles the right-handed-frame correction automatically (no reflections introduced).

    \item \textbf{Permutation resolution.} Define $\sigma$ over atoms of identical atomic number. RMSD is computed as
    \begin{equation}
        \text{RMSD} = \min_{\sigma}\ \sqrt{\frac{1}{N}\,\sum_i \lVert \mathbf{r}_i^{(1)} - \mathbf{R}^* \mathbf{r}_{\sigma(i)}^{(2)} \rVert^2}.
    \end{equation}
    The minimisation strategy depends on problem size:
    \begin{itemize}
        \item If the total number of permutations is $\leq 720 = 6!$, brute-force enumerate all of them and re-run Kabsch for each. This is needed when one element dominates the structure (e.g.\ $\mathrm{HCO_3^-}$ has 3 of 5 atoms as oxygen): the initial Kabsch alignment is then unreliable because most atoms are interchangeable, and Hungarian assignment on a single bad alignment gets stuck in a local minimum.
        \item Otherwise: Hungarian linear-assignment algorithm on the distance matrix from the Kabsch-aligned configuration. For larger molecules ($> 6$ equivalent atoms per element) the alignment from non-equivalent atoms is good enough that Hungarian finds the optimum on the first pass.
    \end{itemize}
\end{enumerate}

\paragraph{Canonical atom ordering.}
Before any comparison the atoms are reordered into the RDKit canonical-rank order. This makes the RMSD deterministic across workers and across different input orderings.

\paragraph{Where used.}
SI-\ref{si:pes} Stage 6 (endpoint distinctness $> \SI{0.25}{\angstrom}$); SI-\ref{si:generative} Stage 6 (endpoint-to-source match $< \SI{1.0}{\angstrom}$); minimum deduplication in the per-compound PES graph (RMSD tiers).


%-----------------------------------------------------------------------------
\subsubsection{Connected-component fragment detection}
\label{sec:shared-fragments}
%-----------------------------------------------------------------------------

A molecular geometry is decomposed into chemically-meaningful fragments by building an adjacency matrix from a covalent-radius-scaled distance criterion,
\begin{equation}
    \text{adj}_{ij} = \mathbb{I}\bigl[\lVert r_i - r_j \rVert < 1.3\,(r^{\text{cov}}_i + r^{\text{cov}}_j)\bigr],
\end{equation}
with covalent radii from ASE. Connected components are extracted by depth-first search and returned as sorted index tensors. The $1.3\times$ scale is conservative; values around $1.15$--$1.20\times$ are more typical for cheminformatics but would mis-classify the stretched bonds at un-converged IRC endpoints as already-dissociated.

Per-fragment validity rules used by both loops: a fragment is accepted iff it is non-empty, its composition satisfies the tetrose cap $\{C{:}5,\ H{:}10,\ O{:}5\}$, and (when single-atom) the atom is hydrogen.

\paragraph{Where used.}
SI-\ref{si:generative} Stage 7 (product-side fragmentation detection); SI-\ref{si:generative} Stage 6 greedy validator (fragmenting both endpoints for chemical-difference comparison); SI-\ref{si:pes} Stage 7 (escape-vs-intramolecular classification by SMILES of relaxed endpoints).


%-----------------------------------------------------------------------------
\subsubsection{SMILES canonicalisation and charge inference}
\label{sec:shared-smiles}
%-----------------------------------------------------------------------------

A 3D geometry is converted to a canonical SMILES via RDKit's bond perception (\texttt{rdkit.Chem.rdDetermineBonds.DetermineBonds}) followed by \texttt{Chem.MolToSmiles(\dots, isomericSmiles=False)}. The non-isomeric form is used everywhere so stereoisomers collapse to the same string --- the diffusion proposer cannot reliably distinguish stereoisomers in any case, and the graph would otherwise admit spurious duplicates that differ only in inverted chirality of a single carbon.

Integer charge is inferred from the perceived bond pattern by summing formal charges atom-by-atom after RDKit's valence-completion pass. The charge is used at three points: (i)~when registering a discovered compound, (ii)~when checking charge conservation across a reaction, and (iii)~when handing a structure to the ML calculator. The DFT refinement (SI-\ref{si:kinetics}) re-derives charges from a wavefunction calculation downstream.


%-----------------------------------------------------------------------------
\subsubsection{Hessian computation and projection}
\label{sec:shared-hessian}
%-----------------------------------------------------------------------------

The Hessian is the central object for both saddle search and validation. The full $3N \times 3N$ matrix is computed by automatic differentiation of the ML forces with respect to atomic positions (\S\ref{sec:shared-mdet}, \texttt{get\_batched\_hessians}), and is then projected to remove the six trivial modes (three translations, three rotations) that contaminate the spectrum at non-zero geometric centring or non-zero overall angular momentum:
\begin{equation}
    \mathbf{H}_\text{proj} = \mathbf{P}\,\mathbf{H}\,\mathbf{P}, \qquad \mathbf{P} = \mathbf{I} - \sum_{k=1}^{6} \mathbf{v}_k \mathbf{v}_k^\top,
\end{equation}
where $\{\mathbf{v}_k\}$ is a mass-weighted orthonormal basis of the translation/rotation subspace obtained by Gram--Schmidt on the six analytical T/R vectors. The projected Hessian is diagonalised into $\{\lambda_i, \mathbf{u}_i\}$ pairs (sorted ascending); eigenvalues with $|\lambda_i| < 10^{-4}\,\text{eV/\AA}^2$ are masked as numerical-noise floor before any chemical interpretation. The two loops apply \emph{different} TS-spectrum policies and we do not unify them: the PES loop (\S\ref{sec:pes-prfo}) admits only \emph{exactly one} significant negative eigenvalue, while the generative loop (\S\ref{sec:hessian-irc}) admits \emph{at least one} below $\SI{-0.01}{eV/\angstrom^2}$. This is a deliberate precision/throughput trade-off: the PES loop trusts its P-RFO walker to converge to a clean first-order saddle, whereas the generative loop accepts higher-order Hessian signatures from diffusion-proposed geometries because the IRC and barrier gates downstream filter out the consequences. Curvature ($\lambda < \SI{-0.01}{eV/\angstrom^2}$) and significance ($|\lambda| < 10^{-4}\,\text{eV/\AA}^2$) thresholds also differ by loop and are stated in the calling SIs.


%-----------------------------------------------------------------------------
\subsubsection{FIRE optimiser}
\label{sec:shared-fire}
%-----------------------------------------------------------------------------

The Fast Inertial Relaxation Engine \citep[][used force-only]{bitzek2006fire} is a velocity-Verlet-based first-order minimiser. Each step:
\begin{enumerate}
    \item Velocity-Verlet update of positions and velocities under the current forces.
    \item Project velocity onto force direction: $\mathbf{v} \leftarrow (1-\alpha)\mathbf{v} + \alpha\,|\mathbf{v}|\,\hat{\mathbf{F}}$.
    \item If $\mathbf{F} \cdot \mathbf{v} > 0$ (descent direction) for the past $N_\text{min}$ steps, increase $\Delta t$ by a factor $f_\text{inc}$ (up to $\Delta t_\text{max}$) and decay $\alpha$ by $f_\alpha$; otherwise reset $\mathbf{v} = 0$, reset $\alpha$ to its initial value, and shrink $\Delta t$ by $f_\text{dec}$.
\end{enumerate}
Default constants follow the original paper: $\Delta t_\text{init} = 0.1$, $\Delta t_\text{max} = 1.0$, $\alpha_\text{start} = 0.1$, $N_\text{min} = 5$. FIRE is the workhorse for any optimisation that does not need a Hessian: per-fragment relax (SI-\ref{si:generative} Stage 7), merge encounter prep (SI-\ref{si:generative} Stage 2), and the NEB backbone (\S\ref{sec:shared-neb}).


%-----------------------------------------------------------------------------
\subsubsection{Trust-region Newton minimiser}
\label{sec:shared-newton}
%-----------------------------------------------------------------------------

When a Hessian is available --- e.g., during IRC endpoint relaxation (SI-\ref{si:pes} Stage 6, SI-\ref{si:generative} Stage 5) --- a second-order trust-region Newton step is used instead of FIRE:
\begin{enumerate}
    \item Project out T/R modes (\S\ref{sec:shared-hessian}), diagonalise, and floor the absolute eigenvalues to $|\lambda_\text{min}| = 10^{-2}\,\text{eV/\AA}^2$ to prevent pathologically large steps in flat (floppy) directions.
    \item Compute the rational-function downhill step in each mode, $p_i = g_i / \lambda_i$, where $g_i$ is the gradient projection on the $i$-th eigenvector.
    \item Constrain the total step to the current trust radius (clamp on per-mode magnitudes if the global norm exceeds the radius).
    \item Evaluate the trial position, compute $\rho = \Delta E_\text{actual} / \Delta E_\text{predicted}$ where $\Delta E_\text{predicted} = -\mathbf{f}\cdot\mathbf{p} + \tfrac{1}{2}\mathbf{p}^\top\mathbf{H}\mathbf{p}$, and accept if $\rho > 0.1$. On rejection, shrink the trust radius by a factor of $2$ and retry (up to 5 retries before giving up).
    \item After acceptance: shrink the trust radius if $\rho < 0.25$ or $\rho > 2.0$ (poor quadratic model), expand by $1.5$ if $0.5 < \rho < 1.5$ and the step hit the trust boundary (good model).
\end{enumerate}
Default trust radius for endpoint relaxation: $\Delta_0 = \SI{0.3}{\angstrom}$, bounded by $[\SI{0.005}{\angstrom}, \SI{2.0}{\angstrom}]$. P-RFO (SI-\ref{si:pes} Stage 4) uses a \emph{partitioned} trust region with a separate TS-mode radius scaled by $1.5\times$ the minimisation radius; that machinery is loop-specific and stays in SI-\ref{si:pes}.


%-----------------------------------------------------------------------------
\subsubsection{Intrinsic Reaction Coordinate (IRC)}
\label{sec:shared-irc}
%-----------------------------------------------------------------------------

Both loops validate a candidate transition state by following the imaginary normal mode in both directions and relaxing each endpoint to a minimum. The shared algorithm:

\begin{enumerate}
    \item Project the Hessian (\S\ref{sec:shared-hessian}); confirm that at least one significant eigenvalue is negative (below the loop-specific curvature threshold; \S\ref{sec:pes-saddle-validation}, \S\ref{sec:hessian-irc}), and take $\hat{\mathbf{u}}_\text{TS}$ to be the eigenvector of the most negative eigenvalue. The IRC step itself does \emph{not} count negative modes beyond this. The strict ``exactly one'' criterion that defines a first-order saddle is enforced \emph{upstream} of the IRC in the PES loop only --- at P-RFO convergence (\S\ref{sec:pes-prfo}), before any candidate reaches this stage --- so by construction the geometries that enter the PES-loop IRC carry a single negative mode. The generative loop has no such upstream guarantee: its candidates come from a diffusion proposer (\S\ref{sec:diffusion}) and may carry $\geq 1$ negative modes at the IRC entry; higher-order saddles are then filtered only implicitly, by failing the endpoint or bidirectional-barrier gates downstream (\S\ref{sec:endpoints}, \S\ref{sec:bwd-barrier}).

    \item Compute a curvature-adaptive displacement magnitude $\Delta x$ from the harmonic-approximation form
    \begin{equation}
        \Delta x = \mathrm{clip}\!\left(\frac{c}{\sqrt{|\lambda_\text{TS}|}},\ \Delta x_\text{min},\ \Delta x_\text{max}\right),
    \end{equation}
    with $c$, $\Delta x_\text{min}$, $\Delta x_\text{max}$ chosen per loop (SI-\ref{si:pes}: $c = \sqrt{0.16}$, range $[0.02, 0.5]\,\si{\angstrom}$; SI-\ref{si:generative}: $c = 0.2$, range $[0.05, 0.5]\,\si{\angstrom}$). Stiff saddles get smaller displacements; flat saddles get larger ones. Curvatures below $10^{-6}\,\text{eV/\AA}^2$ are treated as flat and use the maximum.

    \item Build $\mathbf{R}^\pm = \mathbf{R}_\text{TS} \pm \Delta x\,\hat{\mathbf{u}}_\text{TS}$ and relax each independently to a minimum using trust-region Newton (\S\ref{sec:shared-newton}). Loop-specific convergence: SI-\ref{si:pes} uses $f_\text{max} = \SI{0.005}{eV/\angstrom}$, max $120$ steps; SI-\ref{si:generative} uses the same $f_\text{max}$ but $350$ steps, reflecting the wider distribution of saddle curvatures the diffuser produces. Bail-late acceptance at $f_\text{max} < 10\times$ tolerance is shared.

    \item \emph{Bail-early rule:} if the forward direction fails to relax, skip the backward direction entirely. The saddle is rejected immediately; this saves one expensive relaxation per failure.

    \item Assign direction by energy: the higher-energy endpoint is labelled \emph{reactant} and the lower-energy \emph{product}. This stores every reaction downhill in the ML view, which collapses direction ambiguity for per-pair deduplication. SI-\ref{si:pes} additionally enforces endpoint distinctness (RMSD $> \SI{0.25}{\angstrom}$) and an asymmetric energy-drop criterion; SI-\ref{si:generative} additionally enforces endpoint identity against the source compound (strict + greedy modes).
\end{enumerate}


%-----------------------------------------------------------------------------
\subsubsection{Nudged Elastic Band (NEB)}
\label{sec:shared-neb}
%-----------------------------------------------------------------------------

NEB \citep{henkelman2000climbing} is used \emph{only} for the same-basin same-or-different decision in minimum deduplication: given two relaxed minima with $E_1, E_2$, is there a barrier between them above a threshold?

\paragraph{Band construction.}
The two endpoints are optimally aligned (Kabsch + permutation; \S\ref{sec:shared-rmsd}). The number of interior images is adaptive on the endpoint RMSD,
\begin{equation}
    N_\text{images} = \mathrm{clip}\!\left(\Bigl\lceil \tfrac{\text{RMSD}}{0.15\,\si{\angstrom}} \Bigr\rceil,\ 5,\ 12\right),
\end{equation}
and images are interpolated using the Image-Dependent Pair Potential (IDPP) method via ASE, which produces a physically reasonable initial band even for substantial rearrangements (better than linear interpolation, which routinely creates atom overlaps).

\paragraph{Two-phase optimisation.}
Total FIRE budget $N_\text{FIRE} = 150$ steps with a spring constant $k_\text{spring} = \SI{0.1}{eV/\angstrom^2}$:
\begin{itemize}
    \item \textbf{Phase 1 (regular NEB):} $2/3$ of the budget. All images relaxed onto the MEP with the standard tangent projection of the spring forces.
    \item \textbf{Early bail:} if the peak energy after Phase 1 exceeds $\SI{1.0}{eV}$ above the higher endpoint, the band is considered unphysical (bad interpolation in a system where IDPP could not find a sensible path) and the same-basin check returns ``different basins''.
    \item \textbf{Phase 2 (CI-NEB):} remaining $1/3$ of the budget with the climbing-image variant: the highest-energy image is pushed against the spring forces toward the true saddle, refining the barrier estimate without re-relaxing the rest of the band.
\end{itemize}

\paragraph{Same-basin criterion.}
\begin{equation}
    \Delta E^\ddagger = \max_i(E_i) - \max(E_1, E_2),
\end{equation}
the barrier above the \emph{higher} endpoint. If $\Delta E^\ddagger < \SI{0.05}{eV}$, the two minima are merged. The threshold corresponds to roughly $2\,k_B T$ at room temperature: any barrier smaller than this is freely crossed thermally and treating the two minima as distinct would over-resolve the conformational landscape.


%-----------------------------------------------------------------------------
\subsubsection{Trajectory force integration}
\label{sec:shared-traj}
%-----------------------------------------------------------------------------

When a relaxation trajectory is available, the barrier from the trajectory endpoint to the TS is computed by integrating the force component along the path,
\begin{equation}
    \Delta E^\ddagger \;=\; \int_{\mathbf{R}_\text{min}}^{\mathbf{R}_\text{TS}} \mathbf{F}(\mathbf{r}) \cdot d\mathbf{r}
    \;\approx\; -\,\tfrac{1}{2}\sum_{n} \bigl(\mathbf{F}_n + \mathbf{F}_{n+1}\bigr) \cdot \bigl(\mathbf{R}_{n+1} - \mathbf{R}_{n}\bigr),
\end{equation}
discretised by the trapezoidal rule over the stored trajectory frames (positions and forces, optionally with Hessians for a second-order correction). This is preferred over the endpoint energy difference $E_\text{TS} - E_\text{min}$ because the force integral accumulates over the whole path: spurious local fluctuations in the ML energy (which scale as the per-frame noise of the energy head) average out, while the force head's smoother behaviour dominates the result. When the trajectory is too short for a meaningful integral (or unavailable), the fallback is the endpoint difference.


%-----------------------------------------------------------------------------
\subsubsection{Citation notes}
\label{sec:shared-cites}
%-----------------------------------------------------------------------------

The algorithms above are standard in computational chemistry; this section consolidates them for cross-reference rather than as novel contributions. The primary references implied throughout: Kabsch (1976) for orthogonal Procrustes; Bitzek et al.\ (2006) for FIRE; Banerjee, Adams, Simons \& Shepard (1985) and Baker (1986) for P-RFO (the source-of-truth for the trust-region partition in SI-\ref{si:pes}); Henkelman \& J\'{o}nsson (2000) for the improved-tangent NEB and CI-NEB; Smidstrup et al.\ (2014) for the IDPP interpolation. Loop-specific deviations from these standard implementations (e.g., the $10$-retry overlap guard on encounter-complex displacement in SI-\ref{si:generative}, or the partitioned trust radii for P-RFO) are documented in the respective loop SI.

% --- end si6_shared_machinery.tex ---

% --- begin si4_seeds.tex ---
% ============================================================================
% Supplementary Information SI-4: Seed Compounds and Buffer Equilibria
%
% Drop-in subsection fragment. Parent SI must provide: amsmath/amssymb,
% booktabs, siunitx, mhchem, and the tcolorbox environments {evaluation},
% {concern}, {threshbox}. Bibliography handled at the parent level.
% (Parameter reference lives in si_parameters.tex.)
% ============================================================================

\subsection{Seed Compounds and Buffer Equilibria}
\label{si:seed}

The system is bootstrapped with nine compounds and four hardcoded reactions. The compounds set the chemical scope and provide the unbiased starting conditions for the kinetics ODE; the four reactions inject acid--base equilibria that the diffusion model cannot discover.

\paragraph{Seed compounds.}
Eight molecules are inserted as \texttt{is\_seed = true} compounds at the first worker's startup, under a Postgres advisory lock so only one worker performs the seeding:
\begin{itemize}
    \item $\mathrm{CH_2O}$ (formaldehyde) --- the primary substrate, the only species at non-trace initial concentration in the upstream formose-style runs;
    \item $\mathrm{H_2O}$ (water), $\mathrm{H_2}$ (rendered as \texttt{[HH]} since the MLFF has no electron concept and so $\mathrm{H}$ / $\mathrm{H^+}$ share that SMILES) --- universal small fragments loaded from XYZ files;
    \item $\mathrm{CO_2}$, $\mathrm{H_2CO_3}$, $\mathrm{HCO_3^-}$, $\mathrm{OH^-}$, $\mathrm{H_3O^+}$ --- the buffer species needed to support the manual equilibria.
\end{itemize}
Every seed compound enters with a real 3D geometry (either from \texttt{start.xyz} or the per-fragment \texttt{.xyz} files in \texttt{fragments/}), is canonicalised through RDKit so its SMILES matches what the discovery pipeline would produce, and is registered with a single PES-graph minimum. Larger buffer compounds ($\geq 3$ atoms: $\mathrm{H_2CO_3}$, $\mathrm{HCO_3^-}$) are enqueued for PES exploration and CREST conformer search; smaller ones ($\mathrm{OH^-}$, $\mathrm{H_3O^+}$) have no internal conformers.

Each seed receives an initial concentration of \SI{1}{mM} in the kinetics ODE (SI-\ref{si:kinetics}). The same baseline is applied to every runtime-discovered species so that newly-found compounds can participate in reactions at $t = 0$ instead of being trapped at exactly zero.

\paragraph{Manual buffer equilibria.}
Four acid--base reactions are seeded with literal forward and backward rate constants and tagged \texttt{discovery\_method = "manual\_equilibrium"}:

\begin{center}
\begin{tabular}{lll@{\hspace{2em}}rr}
\toprule
Name & Forward direction & & $k_\text{fwd}$ & $k_\text{bwd}$ \\
\midrule
Water autoionisation
    & $\mathrm{H_2O \to H^+ + OH^-}$
    && $1.8 \times 10^{-3}$ & $10^{10}$ \\
$\mathrm{CO_2}$ hydration
    & $\mathrm{CO_2 + H_2O \to H_2CO_3}$
    && $5.5 \times 10^{-3}$ & $178$ \\
$\mathrm{H_2CO_3}$ dissociation
    & $\mathrm{H_2CO_3 \to H^+ + HCO_3^-}$
    && $2.5 \times 10^{6}$  & $10^{10}$ \\
Proton solvation
    & $\mathrm{H^+ + H_2O \to H_3O^+}$
    && $10^{10}$ & $1.0$ \\
\bottomrule
\end{tabular}
\end{center}

\noindent Units: $\si{s^{-1}}$ for unimolecular forward and \si{M^{-1}.s^{-1}} for bimolecular, by convention; the solver applies the appropriate mass-action exponents from the stoichiometry. The values are literature rate constants from the upstream reference implementation (\texttt{crn-exploration/lib/kinetic\_sampler.py}); they are not derived from Eyring and are \emph{not} re-scaled with temperature when the solver runs at a non-ambient $T$.

The acid--base channels are off-limits to the diffusion proposer for a structural reason: the MD-ET force field used as the energy oracle does not represent explicit electrons, so it has no concept of a proton versus a neutral hydrogen atom. Both render to the same SMILES \texttt{[HH]}, and the model's energy/force predictions for ionic species are systematically wrong in ways that would propagate through the IRC and the barrier filters. Manual equilibria sidestep the issue entirely: their kinetic role (maintaining a self-consistent pH, ensuring $\mathrm{CO_2/H_2CO_3/HCO_3^-}$ buffering, providing a reservoir of solvated protons) is delivered by the literature rate constants, and the diffusion proposer is never asked to discover them.

The rate constants are deliberately fast on the backward direction of the dissociation steps ($10^{10}\,\si{s^{-1}}$, the diffusion-limit cap; SI-\ref{si:kinetics}). The equilibria therefore reach their pH-determined ratios on the shortest timescale resolved by the integrator, so by the time discovered chemistry begins to evolve, the buffer is already at equilibrium and acts as a fast-relaxation reservoir rather than a coupled dynamic player.

% --- end si4_seeds.tex ---

% --- begin si_parameters.tex ---
% ============================================================================
% Supplementary Information: Complete Parameter Reference
%
% Drop-in subsection fragment, shared across SI sections. Parent SI must
% provide: booktabs, siunitx (with \angstrom). No longtable, no tcolorbox.
%
% Organised by SI section, then by pipeline stage. Each stage gets its
% own plain inline tabular so that LaTeX can break naturally between
% stages across pages.
% ============================================================================

\subsection{Complete Parameter Reference}
\label{si:parameters}

This subsection consolidates every numerical parameter, threshold, and tolerance used by the discovery pipelines and their supporting machinery, at the values used in the main run.

\bigskip

%-----------------------------------------------------------------------------
\subsubsection{PES loop (SI-\ref{si:pes})}
\label{sec:params-pes}
%-----------------------------------------------------------------------------

\noindent\textbf{Stage 2 --- Molecular dynamics.}\par\nopagebreak
\noindent\begin{tabular}{p{5cm}rl}
\toprule
Parameter & Value & Unit \\
\midrule
MD steps & 2500 & --- \\
Integration step & 0.5 & fs \\
Total simulation time & 1.25 & ps \\
Temperature & 300 & K \\
Thermostat & Langevin & --- \\
Coupling $\tau$ & 100 & fs \\
Snapshot interval & 10 & steps \\
\bottomrule
\end{tabular}

\bigskip

\noindent\textbf{Stage 3 --- TS candidates.}\par\nopagebreak
\noindent\begin{tabular}{p{5cm}rl}
\toprule
Parameter & Value & Unit \\
\midrule
Candidates per iteration & 32 & --- \\
Selection & uniform random & --- \\
\bottomrule
\end{tabular}

\bigskip

\noindent\textbf{Stage 4 --- P-RFO.}\par\nopagebreak
\noindent\begin{tabular}{p{5cm}rl}
\toprule
Parameter & Value & Unit \\
\midrule
Max steps & 120 & --- \\
$f_\text{max}$ tolerance & 0.01 & eV/\AA \\
$f_\text{RMS}$ tolerance & 0.005 & eV/\AA \\
Initial trust radius & 0.5 & \AA \\
Trust radius bounds & $[0.01, 0.3]$ & \AA \\
TS-mode trust scale & 1.5 & --- \\
Accept threshold ($\rho$) & 0.1 & --- \\
Max retries / step & 5 & --- \\
Early-stop window & 20 & steps \\
Early-stop displacement & $10^{-4}$ & \AA \\
Early-stop efficiency & 0.1 & --- \\
\bottomrule
\end{tabular}

\bigskip

\noindent\textbf{Stage 5 --- Saddle validation.}\par\nopagebreak
\noindent\begin{tabular}{p{5cm}rl}
\toprule
Parameter & Value & Unit \\
\midrule
Significance threshold & $10^{-4}$ & eV/\AA$^2$ \\
TS eigenvalue threshold & $-0.01$ & eV/\AA$^2$ \\
Required spectrum & exactly 1 neg & --- \\
\bottomrule
\end{tabular}

\bigskip

\noindent\textbf{Stage 6 --- Bidirectional IRC.}\par\nopagebreak
\noindent\begin{tabular}{p{5cm}rl}
\toprule
Parameter & Value & Unit \\
\midrule
Energy-drop target & 0.08 & eV \\
Displacement formula & $\sqrt{2 \cdot 0.08 / |\lambda|}$ & \AA \\
Displacement clamp & $[0.02, 0.5]$ & \AA \\
Endpoint $f_\text{max}$ & 0.005 & eV/\AA \\
Endpoint max steps & 120 & --- \\
Energy threshold (symmetric) & 0.04 & eV \\
Energy threshold (asymmetric) & 0.005 & eV \\
Min endpoint RMSD & 0.25 & \AA \\
\bottomrule
\end{tabular}

\bigskip

\noindent\textbf{Stage 7 --- Classification.}\par\nopagebreak
\noindent\begin{tabular}{p{5cm}rl}
\toprule
Parameter & Value & Unit \\
\midrule
Source-match test & canonical SMILES & --- \\
Categories & intramol / escape / anomalous & --- \\
\bottomrule
\end{tabular}

\bigskip

\noindent\textbf{Stage 8 --- Commitment.}\par\nopagebreak
\noindent\begin{tabular}{p{5cm}rl}
\toprule
Parameter & Value & Unit \\
\midrule
Intramol minimum dedup tiers & 3 & --- \\
Ultra-low merge & 0.01 & \AA{} RMSD \\
Fast merge & 0.20, 0.05 & \AA{}, eV \\
NEB same-basin threshold & 0.05 & eV barrier \\
NEB images, FIRE steps & 5--12, 150 & --- \\
TS barrier-match tolerance & 0.1 & eV \\
Escape pipeline & SI-\ref{si:generative} 4--8 & --- \\
Barrier preference & traj.\ force integral & --- \\
\bottomrule
\end{tabular}

\bigskip

%-----------------------------------------------------------------------------
\subsubsection{Generative loop (SI-\ref{si:generative})}
\label{sec:params-gen}
%-----------------------------------------------------------------------------

\noindent\textbf{Stage 1 --- Sampling.}\par\nopagebreak
\noindent\begin{tabular}{p{5cm}rl}
\toprule
Parameter & Value & Unit \\
\midrule
Batch size & 32 & contexts \\
Single/merge split & 0.5 & --- \\
Oversample factor (merges) & 50 & --- \\
Min compound size & 3 & atoms \\
Composition cap & C5H10O5 & --- \\
Per-pair reaction cap & 4 & --- \\
Boltzmann $k_B T$ & 0.0257 & eV \\
vdW separation buffer & 0.2 & \AA \\
\bottomrule
\end{tabular}

\bigskip

\noindent\textbf{Stage 2 --- Merge enhancement.}\par\nopagebreak
\noindent\begin{tabular}{p{5cm}rl}
\toprule
Parameter & Value & Unit \\
\midrule
Extra orientations ($K$) & 5 & --- \\
FIRE micro-relax steps & 5 & --- \\
Symmetry-break $\sigma$ & 0.15 & \AA \\
Encounter-complex cap & 0.3 & eV \\
Min pairwise distance & 0.8 & \AA \\
\bottomrule
\end{tabular}

\bigskip

\noindent\textbf{Stage 3 --- Diffusion.}\par\nopagebreak
\noindent\begin{tabular}{p{5cm}rl}
\toprule
Parameter & Value & Unit \\
\midrule
$T$ (schedule length) & 1000 & steps \\
Polynomial power $p$ & 2 & --- \\
Precision $s$ & $10^{-5}$ & --- \\
Clip floor & $10^{-3}$ & --- \\
Partial-noise range & 100, 650 & steps \\
Max reverse steps & 1000 & --- \\
Convergence step & 0 & --- \\
Cutoff (neighbour list) & 5.0 & \AA \\
\bottomrule
\end{tabular}

\bigskip

\noindent\textbf{Stage 4 --- Forward barrier.}\par\nopagebreak
\noindent\begin{tabular}{p{5cm}rl}
\toprule
Parameter & Value & Unit \\
\midrule
Acceptance band & $[0,\ 2.7211]$ & eV \\
\bottomrule
\end{tabular}

\bigskip

\noindent\textbf{Stage 5 --- Hessian + IRC.}\par\nopagebreak
\noindent\begin{tabular}{p{5cm}rl}
\toprule
Parameter & Value & Unit \\
\midrule
TS eigenvalue threshold & $-0.01$ & eV/\AA$^2$ \\
Significance mask & $10^{-4}$ & eV/\AA$^2$ \\
Required spectrum & $\geq 1$ neg & --- \\
Adaptive step formula & $0.2 / \sqrt{|\lambda|}$ & \AA \\
Step clamp & $[0.05, 0.5]$ & \AA \\
IRC max steps / direction & 350 & --- \\
IRC $f_\text{max}$ & 0.005 & eV/\AA \\
Bail-late tolerance & $10\times f_\text{max}$ & --- \\
\bottomrule
\end{tabular}

\bigskip

\noindent\textbf{Stage 6 --- Endpoint identity.}\par\nopagebreak
\noindent\begin{tabular}{p{5cm}rl}
\toprule
Parameter & Value & Unit \\
\midrule
RMSD match threshold & 1.0 & \AA \\
Greedy fallback & enabled & --- \\
Max fragments / side & 2 & --- \\
\bottomrule
\end{tabular}

\bigskip

\noindent\textbf{Stage 7 --- Fragmentation.}\par\nopagebreak
\noindent\begin{tabular}{p{5cm}rl}
\toprule
Parameter & Value & Unit \\
\midrule
Covalent-radius scale & 1.3 & --- \\
Per-fragment composition cap & C5H10O5 & --- \\
Per-fragment $f_\text{max}$ & 0.005 & eV/\AA \\
Per-fragment max steps & 200 & --- \\
\bottomrule
\end{tabular}

\bigskip

\noindent\textbf{Stage 8 --- Backward barrier.}\par\nopagebreak
\noindent\begin{tabular}{p{5cm}rl}
\toprule
Parameter & Value & Unit \\
\midrule
Acceptance band & $(0,\ 2.7211]$ & eV \\
Barrier computation & traj.\ force integral & --- \\
\bottomrule
\end{tabular}


%-----------------------------------------------------------------------------
\subsubsection{Kinetics solver (SI-\ref{si:kinetics})}
\label{sec:params-kin}
%-----------------------------------------------------------------------------

\noindent\textbf{Rate constants.}\par\nopagebreak
\noindent\begin{tabular}{p{5cm}rl}
\toprule
Parameter & Value & Unit \\
\midrule
Boltzmann $k_B$ & $8.617 \times 10^{-5}$ & eV/K \\
Planck $h$ & $4.136 \times 10^{-15}$ & eV\,s \\
Diffusion limit $k_\text{diff}$ & $10^{10}$ & M$^{-1}$s$^{-1}$ \\
Default exploration $T$ & 500 & K \\
Barrier overflow cap & 10 & eV \\
Barrier-priority order & DFT $>$ ML in-box $>$ drop & --- \\
\bottomrule
\end{tabular}

\bigskip

\noindent\textbf{Build-time filters.}\par\nopagebreak
\noindent\begin{tabular}{p{5cm}rl}
\toprule
Parameter & Value & Unit \\
\midrule
Disconnected-SMILES rule & drop & --- \\
Negative-$E_a$ tolerance & $-0.05$ & eV \\
Barrier cutoff & 10 & eV \\
Direction-agnostic dedup & lowest $\min(E_a^\text{fwd}, E_a^\text{bwd})$ & --- \\
\bottomrule
\end{tabular}

\bigskip

\noindent\textbf{Species universe and initial conditions.}\par\nopagebreak
\noindent\begin{tabular}{p{5cm}rl}
\toprule
Parameter & Value & Unit \\
\midrule
Uniform baseline & $10^{-3}$ & M \\
Seed species ($9$) & $10^{-3}$ & M \\
Noise floor & $10^{-20}$ & M \\
\bottomrule
\end{tabular}

\bigskip

\noindent\textbf{Solver (PETSc BDF).}\par\nopagebreak
\noindent\begin{tabular}{p{5cm}rl}
\toprule
Parameter & Value & Unit \\
\midrule
$t_\text{max}$ & $10^{8}$ & s \\
Initial step $\Delta t_0$ & $10^{-10}$ & s \\
Absolute tolerance & $10^{-16}$ & --- \\
Relative tolerance & $10^{-12}$ & --- \\
Max steps & $5 \times 10^{7}$ & --- \\
Newton linear solver & preonly + LU & --- \\
Output grid (plot) & 400 points, $[10^{-12}, 10^{8}]$~s & geometric \\
Output grid (decade) & 19 points, $10^{-10}$--$10^{8}$~s & --- \\
\bottomrule
\end{tabular}

\bigskip

\noindent\textbf{Polling and trigger.}\par\nopagebreak
\noindent\begin{tabular}{p{5cm}rl}
\toprule
Parameter & Value & Unit \\
\midrule
Poll interval & 60 & s \\
Network-version threshold & 5 & reactions \\
Max snapshot age & 600 & s \\
DFT-delta trigger & $> 0$ new & --- \\
Per-experiment advisory locks & key 43, 44, $\ldots$ & --- \\
\bottomrule
\end{tabular}

\bigskip

\noindent\textbf{Steady-state distribution (sampler input).}\par\nopagebreak
\noindent\begin{tabular}{p{5cm}rl}
\toprule
Parameter & Value & Unit \\
\midrule
Form & softmax$(\log_{10} c)$ at $t_\text{max}$ & --- \\
Exclusion rule & $c \leq 10^{-20}$~M dropped & --- \\
Min species for valid snapshot & 2 & --- \\
\bottomrule
\end{tabular}

\bigskip

\noindent\textbf{DFT energy refinement (\S\ref{sec:kinetics-dft}).}\par\nopagebreak
\noindent\begin{tabular}{p{5cm}rl}
\toprule
Parameter & Value & Unit \\
\midrule
Functional & PBE0 & --- \\
Basis set & def2-TZVPP & --- \\
SCF convergence (\texttt{conv\_tol}) & $10^{-9}$ & Hartree \\
SCF max cycles & 300 & --- \\
Fallback on non-convergence & Newton solver & --- \\
Spin handling & singlet/doublet from $N_e \bmod 2$ & --- \\
Per-reaction calculations & 3 single-points (R, TS, P) & --- \\
Per-compound calculation & 1 single-point at lowest-$E$ min & --- \\
Hessians / frequency check & not computed & --- \\
Geometry re-optimisation & not performed & --- \\
Typical wall time / TS & $\sim 10$~s to $\sim 5$~min & --- \\
\bottomrule
\end{tabular}


%-----------------------------------------------------------------------------
\subsubsection{Seed compounds and buffer equilibria (SI-\ref{si:seed})}
\label{sec:params-seed}
%-----------------------------------------------------------------------------

\noindent\textbf{Seed compounds.}\par\nopagebreak
\noindent\begin{tabular}{p{4cm}p{4cm}rp{4cm}}
\toprule
SMILES & Name & Charge \\
\midrule
\texttt{C=O}        & formaldehyde   & $0$ \\
\texttt{[HH]}       & H / H$^+$      & $0$ \\
\texttt{O}          & water          & $0$ \\
\texttt{[OH-]}      & hydroxide      & $-1$ \\
\texttt{[OH3+]}     & hydronium      & $+1$ \\
\texttt{O=C=O}      & CO$_2$         & $0$ \\
\texttt{OC(O)O}     & H$_2$CO$_3$    & $0$ \\
\texttt{O=C([O-])O} & HCO$_3^-$      & $-1$ \\
\bottomrule
\end{tabular}

\bigskip

\noindent\textbf{Manual buffer-equilibrium reactions.}\par\nopagebreak
\noindent\begin{tabular}{p{4cm}p{5cm}rr}
\toprule
Name & Forward & $k_\text{fwd}$ \\
\midrule
Water autoionisation        & \texttt{O -> [HH] + [OH-]}              & $1.8 \times 10^{-3}$ \\
CO$_2$ hydration            & \texttt{O=C=O + O -> OC(O)O}            & $5.5 \times 10^{-3}$ \\
H$_2$CO$_3$ dissociation    & \texttt{OC(O)O -> [HH] + O=C([O-])O}    & $2.5 \times 10^{6}$ \\
Proton solvation            & \texttt{[HH] + O -> [OH3+]}             & $10^{10}$ \\
\bottomrule
\end{tabular}
\\[1ex]
\noindent\textit{Units: $\si{s^{-1}}$ (unimolecular), $\si{M^{-1}\cdot s^{-1}}$ (bimolecular). Rate constants bypass Eyring and are temperature-independent.}


%-----------------------------------------------------------------------------
\subsubsection{Shared machinery: ML force field, geometry, optimisation (SI-\ref{si:shared})}
\label{sec:params-shared}
%-----------------------------------------------------------------------------

\noindent\textbf{ML force field (\S\ref{sec:shared-mdet}).}\par\nopagebreak
\noindent\begin{tabular}{p{5cm}rl}
\toprule
Parameter & Value & Unit \\
\midrule
Model & MD-ET transformer (direct force) & --- \\
Variant & \texttt{12l} (12 attention layers) & --- \\
Cutoff (neighbour list) & 5.0 & \AA \\
Training data & QCML & --- \\
Hessian source & autograd of forces & --- \\
Hessian batch size & 8 & structures \\
\bottomrule
\end{tabular}

\bigskip

\noindent\textbf{RMSD (\S\ref{sec:shared-rmsd}).}\par\nopagebreak
\noindent\begin{tabular}{p{5cm}rl}
\toprule
Parameter & Value & Unit \\
\midrule
Alignment & Kabsch + SVD & --- \\
Permutation (small) & brute-force, $\leq 720$ perms & --- \\
Permutation (large) & Hungarian assignment & --- \\
Canonical atom ordering & RDKit canonical rank & --- \\
\bottomrule
\end{tabular}

\bigskip

\noindent\textbf{Fragment detection (\S\ref{sec:shared-fragments}).}\par\nopagebreak
\noindent\begin{tabular}{p{5cm}rl}
\toprule
Parameter & Value & Unit \\
\midrule
Covalent-radius scale & 1.3 & --- \\
Validity: non-empty & --- & --- \\
Validity: single-atom rule & H only & --- \\
Validity: composition cap & $\{C{:}5, H{:}10, O{:}5\}$ & --- \\
\bottomrule
\end{tabular}

\bigskip

\noindent\textbf{Hessian + T/R projection (\S\ref{sec:shared-hessian}).}\par\nopagebreak
\noindent\begin{tabular}{p{5cm}rl}
\toprule
Parameter & Value & Unit \\
\midrule
Projection basis & 6 mass-weighted T/R modes & --- \\
Significance mask & $10^{-4}$ & eV/\AA$^2$ \\
\bottomrule
\end{tabular}

\bigskip

\noindent\textbf{FIRE optimiser (\S\ref{sec:shared-fire}).}\par\nopagebreak
\noindent\begin{tabular}{p{5cm}rl}
\toprule
Parameter & Value & Unit \\
\midrule
$\Delta t_\text{init}$, $\Delta t_\text{max}$ & 0.1, 1.0 & --- \\
$\alpha_\text{start}$ & 0.1 & --- \\
$N_\text{min}$ & 5 & steps \\
\bottomrule
\end{tabular}

\bigskip

\noindent\textbf{Trust-region Newton (\S\ref{sec:shared-newton}).}\par\nopagebreak
\noindent\begin{tabular}{p{5cm}rl}
\toprule
Parameter & Value & Unit \\
\midrule
Eigenvalue floor & $|\lambda_\text{min}| = 10^{-2}$ & eV/\AA$^2$ \\
Initial trust radius & 0.3 & \AA \\
Trust radius bounds & $[0.005, 2.0]$ & \AA \\
Accept threshold ($\rho$) & 0.1 & --- \\
Shrink on $\rho < 0.25$ or $\rho > 2$ & $\times 0.5$ & --- \\
Expand on $0.5 < \rho < 1.5$ \& boundary hit & $\times 1.5$ & --- & Good fit \\
Max retries per step & 5 & --- \\
\bottomrule
\end{tabular}

\bigskip

\noindent\textbf{IRC (\S\ref{sec:shared-irc}).}\par\nopagebreak
\noindent\begin{tabular}{p{5cm}rl}
\toprule
Parameter & Value & Unit \\
\midrule
Displacement form & $c / \sqrt{|\lambda_\text{TS}|}$ & \AA \\
Direction assignment & higher energy $=$ reactant & --- \\
Bail-early rule & forward fail $\to$ skip backward & --- \\
Bail-late tolerance & $f_\text{max} < 10\times$ convergence & --- \\
\bottomrule
\end{tabular}

\bigskip

\noindent\textbf{NEB same-basin check (\S\ref{sec:shared-neb}).}\par\nopagebreak
\noindent\begin{tabular}{p{5cm}rl}
\toprule
Parameter & Value & Unit \\
\midrule
Image count & $\lceil \text{RMSD}/0.15 \rceil$ clip $[5, 12]$ & --- \\
Initial interpolation & IDPP & --- \\
Spring constant $k$ & 0.1 & eV/\AA$^2$ \\
Total FIRE steps & 150 & --- \\
Phase-1 early bail & $\Delta E^\ddagger > 1.0$ eV & --- \\
Same-basin threshold & 0.05 & eV \\
\bottomrule
\end{tabular}

\bigskip

\noindent\textbf{Trajectory force-integral barrier (\S\ref{sec:shared-traj}).}\par\nopagebreak
\noindent\begin{tabular}{p{5cm}rl}
\toprule
Parameter & Value & Unit \\
\midrule
Integration rule & trapezoidal & --- \\
Fallback & endpoint energy difference & --- \\
\bottomrule
\end{tabular}

% --- end si_parameters.tex ---

